\documentclass[11pt]{article}

%%% Packages
\usepackage{amsmath,amssymb,amsthm}
\usepackage{mathtools}
\usepackage{enumerate}
\usepackage[shortlabels]{enumitem}
\usepackage{booktabs}
\usepackage{hyperref}
\usepackage[all]{xy}

%%% Page geometry
\usepackage[margin=1.15in]{geometry}

%%% Theorem environments
\theoremstyle{plain}
\newtheorem{theorem}{Theorem}[section]
\newtheorem{proposition}[theorem]{Proposition}
\newtheorem{corollary}[theorem]{Corollary}
\newtheorem{lemma}[theorem]{Lemma}
\newtheorem{conjecture}[theorem]{Conjecture}

\theoremstyle{definition}
\newtheorem{definition}[theorem]{Definition}
\newtheorem{construction}[theorem]{Construction}
\newtheorem{example}[theorem]{Example}

\theoremstyle{remark}
\newtheorem{remark}[theorem]{Remark}

%%% Macros
\newcommand{\fg}{\mathfrak{g}}
\newcommand{\fh}{\mathfrak{h}}
\newcommand{\Defcyc}{\mathrm{Def}_{\mathrm{cyc}}}
\newcommand{\Ran}{\mathrm{Ran}}
\newcommand{\Mbar}{\overline{\mathcal{M}}}
\newcommand{\Sh}{\mathrm{Sh}}
\newcommand{\FM}{\mathrm{FM}}
\newcommand{\MC}{\mathrm{MC}}
\newcommand{\cA}{\mathcal{A}}
\newcommand{\cH}{\mathcal{H}}
\newcommand{\cW}{\mathcal{W}}
\newcommand{\Res}{\operatorname{Res}}
\newcommand{\disc}{\operatorname{disc}}
\DeclareMathOperator{\Aut}{Aut}

%%% Title
\title{Shadow obstruction towers and the algebraicity of chiral deformation invariants}
\author{Raeez Lorgat}
\date{}

\begin{document}
\maketitle

%%% ================================================================
%%%  ABSTRACT
%%% ================================================================

\begin{abstract}
For a chirally Koszul vertex algebra~$A$ with modular
characteristic~$\kappa$, cubic shadow~$\alpha$, and quartic
shadow~$S_4$, we prove that the infinite sequence of higher
deformation invariants $\{S_r\}_{r \geq 2}$ extracted from the
bar construction is algebraic of degree~$2$:\ the weighted
generating function $H(t) = \sum_{r \geq 2} r\,S_r\,t^r$
satisfies $H(t)^2 = t^4\,Q(t)$, where
\[
  Q(t) \;=\; 4\kappa^2 + 12\kappa\alpha\,t
  + (9\alpha^2 + 16\kappa S_4)\,t^2
\]
is the \emph{shadow metric}, a quadratic polynomial determined
by the genus-$0$ OPE data of~$A$.
The critical discriminant $\Delta = 8\kappa S_4$ partitions all
chirally Koszul algebras into four classes:
Gaussian ($\Delta = 0$, $\alpha = 0$; depth~$2$),
Lie ($\Delta = 0$, $\alpha \neq 0$; depth~$3$),
contact (depth~$4$, by stratum separation), and
mixed ($\Delta \neq 0$; depth~$\infty$).
The entire infinite tail $\{S_r\}_{r \geq 5}$ is determined by the
three invariants $(\kappa, \alpha, S_4)$.
We compute shadow obstruction towers explicitly for all standard families
(Heisenberg, affine Kac--Moody, $\beta\gamma$, Virasoro,
$\cW_N$), derive the shadow connection and its Koszul
monodromy~$-1$, and extract the genus-$2$ cross-channel
correction $\delta F_2(\cW_3) = (c+204)/(16c)$
(the multi-generator universality problem is resolved
negatively: this correction is nonvanishing; see
Remark~\ref{rem:cross-channel-status}) by a
graph-by-graph stable graph sum over the seven genus-$2$ stable
graphs.
\end{abstract}

%%% ================================================================
%%%  1. INTRODUCTION
%%% ================================================================

\section{Introduction}\label{sec:intro}

\subsection{Vertex algebras and moduli of curves}

A vertex algebra $A$ on a smooth projective curve~$X$ produces, for
each genus $g \geq 0$ and each integer $n \geq 0$, a space of
conformal blocks $\mathbb{V}_{g,n}(A)$ on the moduli space
$\Mbar_{g,n}$ of stable curves.  When $A$ is chirally
Koszul (Definition~\ref{def:chirally-koszul}), these assemble
into cohomological data on $\Mbar_{g,n}$, and the problem of
computing them reduces to the operator product expansion (OPE).

The classical approach to this computation, through sewing and
factorization, expresses genus-$g$ amplitudes as sums over
stable graphs $\Gamma$ of genus~$g$:\ each vertex of $\Gamma$
contributes a genus-$g_v$ conformal block, and each edge
contributes a propagator.  The resulting graph sums are the
Feynman expansion of a two-dimensional conformal field theory,
and they produce characteristic classes on $\Mbar_{g,n}$ whose
integrals are the free energies~$F_g(A)$.

We prove that, for any chirally Koszul vertex algebra, these
free energies and the full tower of deformation invariants that
underlies them are controlled by a single algebraic relation
of degree~$2$.


\subsection{The bar construction and the shadow obstruction tower}

The algebraic framework is the chiral bar construction, a
vertex-algebraic analogue of the classical bar construction in
homological algebra.  For a vertex algebra~$A$ on a curve~$X$,
the bar complex $B(A)$ is a factorization coalgebra on the Ran
space $\Ran(X)$, whose differential encodes collision residues
of the OPE on successive Fulton--MacPherson compactifications
$\FM_n(X)$.  When $A$ is \emph{chirally Koszul}, so that
the bar cohomology $H^{\bullet}(B(A))$ is concentrated in bar
degree~$1$, the bar complex carries a Maurer--Cartan element
\[
  \Theta_A \;:=\; D_A - d_0
  \;\in\; \MC\bigl(\Defcyc^{\mathrm{mod}}(A)\bigr),
\]
where $D_A$ is the full bar differential (incorporating
contributions from all genera and all boundary strata of
$\Mbar_{g,n}$) and $d_0$ is the genus-$0$ linear part.
The equation $D_A^2 = 0$, which is a formal consequence of the
boundary-of-boundary relation $\partial^2 = 0$ on $\Mbar_{g,n}$,
ensures that $\Theta_A$ satisfies the Maurer--Cartan equation
\begin{equation}\label{eq:intro-mc}
  d_0\,\Theta_A + \tfrac{1}{2}[\Theta_A, \Theta_A] = 0
\end{equation}
in the modular convolution algebra
$\Defcyc^{\mathrm{mod}}(A)$.  This single equation packages
all genera and all arities simultaneously.

The \emph{shadow obstruction tower} of~$A$ is the sequence of
finite-order projections of~$\Theta_A$:
\[
  \Theta_A^{\leq 2},\quad
  \Theta_A^{\leq 3},\quad
  \Theta_A^{\leq 4},\quad
  \ldots
\]
where $\Theta_A^{\leq r}$ retains only the components of
arity at most~$r$ in the convolution algebra.
On a one-dimensional primary line $L \subset
\Defcyc^{\mathrm{mod}}(A)$ spanned by a single generator~$\eta$
of conformal weight~$h$, the arity-$r$ projection is a scalar
$S_r = S_r(A) \in \mathbb{Q}(c)$:
\begin{itemize}
\item $S_2 = \kappa(A)$, the \emph{modular characteristic}:
  the genus-$1$ curvature scalar determining $F_1 = \kappa/24$.
\item $S_3 = \alpha(A)$, the \emph{cubic shadow}:
  the three-point collision residue on $\FM_3(X)$.
\item $S_4 = S_4(A)$, the \emph{quartic shadow}:
  the four-point collision residue, encoding the first genuinely
  nonlinear OPE invariant.
\item $S_r$ for $r \geq 5$:\ the higher shadows.
\end{itemize}
The main question is: \emph{how much of the OPE is needed to
determine the full tower?}


\subsection{Main results}

Our principal theorem asserts that the answer is \emph{three
invariants}:\ the entire infinite tower $\{S_r\}_{r \geq 2}$ is
algebraic of degree~$2$, determined by the genus-$0$ triple
$(\kappa, \alpha, S_4)$.

\begin{theorem}[Riccati algebraicity; Theorem~\ref{thm:riccati}]
\label{thm:intro-riccati}
Let $A$ be a chirally Koszul vertex algebra, $L$ a primary line
with shadow data $S_r$, and $\kappa = S_2 \neq 0$.  Define
\[
  Q(t)
  \;=\; 4\kappa^2 + 12\kappa\alpha\,t
  + (9\alpha^2 + 16\kappa S_4)\,t^2
\]
and $H(t) = \sum_{r \geq 2} r\,S_r\,t^r$.  Then
\begin{equation}\label{eq:intro-algebraicity}
  H(t)^2 \;=\; t^4\,Q(t).
\end{equation}
Every shadow coefficient $S_r$ with $r \geq 5$ is determined by
$(\kappa, \alpha, S_4)$ via the closed form
$S_r = \frac{1}{r}[t^{r-2}]\sqrt{Q(t)}$.
\end{theorem}

The algebraicity theorem has an immediate structural consequence.
Define the \emph{critical discriminant}
$\Delta := 8\kappa S_4$, so that $Q$ admits the
Gaussian decomposition
\[
  Q(t) \;=\; (2\kappa + 3\alpha\,t)^2 + 2\Delta\,t^2.
\]
The vanishing or nonvanishing of $\Delta$ controls the depth
of the shadow obstruction tower:

\begin{theorem}[Four-class partition; Theorem~\ref{thm:dichotomy}]
\label{thm:intro-classification}
On any primary line $L$ with $\kappa \neq 0$, the shadow
depth $r_{\max} := \sup\{r : S_r \neq 0\}$ belongs to
$\{2, 3, \infty\}$, classified as follows.
\begin{enumerate}[(i)]
\item \textbf{Class G} (Gaussian):
  $\Delta = 0$, $\alpha = 0$.
  $r_{\max} = 2$.  Example:\ Heisenberg.
\item \textbf{Class L} (Lie):
  $\Delta = 0$, $\alpha \neq 0$.
  $r_{\max} = 3$.  Example:\ affine Kac--Moody.
\item \textbf{Class M} (mixed):
  $\Delta \neq 0$.
  $r_{\max} = \infty$, and $S_r = \Delta \cdot R_r$ for
  universal rational functions
  $R_r \in \mathbb{Q}(\alpha, \kappa, S_4)$.
  Example:\ Virasoro, $\cW_N$.
\end{enumerate}
A fourth class, \textbf{class C} (contact, $r_{\max} = 4$),
arises by stratum separation: the quartic contact invariant
lives on a charged stratum that exits the single-line complex.
Example:\ $\beta\gamma$.
\end{theorem}

The shadow obstruction tower also carries a logarithmic connection.

\begin{theorem}[Shadow connection; Theorem~\ref{thm:connection}]
\label{thm:intro-connection}
The shadow metric $Q(t)$ defines a logarithmic connection
$\nabla^{\mathrm{sh}} = d - \frac{Q'(t)}{2Q(t)}\,dt$
on $L \setminus \{Q = 0\}$, with regular singularities
of residue~$\tfrac{1}{2}$ at the zeros of~$Q$.
The monodromy around each branch point is~$-1$.
Under Koszul duality $A \mapsto A^{!}$, the connection transforms
by the Verdier involution; for Virasoro:
\[
  \Delta(\mathrm{Vir}_c)
  + \Delta(\mathrm{Vir}_{26-c})
  \;=\; \frac{6960}{(5c+22)(152-5c)}.
\]
\end{theorem}

Finally, the algebraicity theorem yields explicit genus-$2$
free energies by stable graph summation.

\begin{theorem}[Genus-$2$ universality; Theorem~\ref{thm:genus2}]
\label{thm:intro-genus2}
For any chirally Koszul vertex algebra~$A$ with a single
generator:
\[
  F_2(A) \;=\; \kappa(A) \cdot \lambda_2^{\mathrm{FP}},
  \qquad
  \lambda_2^{\mathrm{FP}} \;=\; \frac{7}{5760}.
\]
For the $\cW_3$ algebra with two generators
$T$ (weight~$2$) and $W$ (weight~$3$), and conditional
on multi-generator universality at genus~$\geq 2$
\textup{(}Remark~\textup{\ref{rem:cross-channel-status})},
the multi-channel
stable graph sum over the seven genus-$2$ stable graphs
yields a scalar part
$F_2^{\mathrm{scal}} = (5c/6) \cdot 7/5760 = 7c/6912$
and a cross-channel correction
\[
  \delta F_2(\cW_3) \;=\;
  \frac{c + 204}{16c},
\]
with per-graph contributions $3/c$ \textup{(}banana\textup{)},
$9/(2c)$ \textup{(}theta\textup{)},
$1/16$ \textup{(}lollipop\textup{)}, and
$21/(4c)$ \textup{(}barbell\textup{)}.
The lollipop contribution is independent of~$c$.
\end{theorem}


\subsection{Context}

The algebraicity relation $H(t)^2 = t^4 Q(t)$ is structurally
parallel to the loop equations of matrix models
\cite{ADKMV,DijkgraafVafa} and to the Eynard--Orantin
topological recursion \cite{EO}, which compute higher-genus
invariants from a spectral curve.  In those frameworks the
spectral curve is an input; here the spectral
curve $\Sigma = \{H^2 = t^4 Q(t)\}$ is an \emph{output} of
the algebraicity theorem, determined by the vertex algebra~$A$.

The shadow connection $\nabla^{\mathrm{sh}}$ has Stokes data
controlled by the critical discriminant: monodromy is trivial
for classes G and~L ($\Delta = 0$) and nontrivial for class~M
($\Delta \neq 0$), in parallel with the wall-crossing
formalism of Kontsevich--Soibelman
\cite{KontsevichSoibelman}.  The invariants live in
the cyclic deformation complex of a vertex algebra, not in
a motivic Hall algebra.

In classical Koszul duality, a Koszul algebra is determined by
its quadratic dual~$A^{!}$, and the bar resolution concentrates
in degree~$1$.  The chiral bar construction plays the same role,
but the modular extension introduces genuinely new data: the
shadow invariants $S_r$ measure the complexity of higher-genus,
higher-arity components of the bar differential, which have no
classical analogue.  The four-class partition
(Theorem~\ref{thm:intro-classification}) classifies this
complexity within the Koszul world; all four classes are Koszul.


\subsection{Notation and conventions}

Throughout, $X$ denotes a smooth projective algebraic curve
over a field~$k$ of characteristic zero.  We work in the
cohomological convention: all differentials have degree~$+1$.
The bar construction uses desuspension ($s^{-1}$, degree~$-1$).

A \emph{vertex algebra} is a vertex algebra in the sense of
\cite{FrenkelBenZvi,Kac-vertex}, equipped with a conformal
vector~$\omega$ (the Virasoro element), a conformal grading
$A = \bigoplus_{h \geq 0} A_h$ with $\dim A_h < \infty$, and
the operator product expansion
$a(z)b(w) \sim \sum_{n \geq 0} (a_{(n)} b)(w) \cdot
(z-w)^{-n-1}$.

The \emph{modular characteristic} $\kappa(A) \in \mathbb{Q}(c)$
is the genus-$1$ curvature scalar: the coefficient in
$F_1(A) = \kappa(A)/24$.  The \emph{central charge}~$c$ is
the standard Virasoro central charge.  These are distinct
invariants:\ $\kappa = k$ for the rank-$1$ Heisenberg at level~$k$,
$\kappa = c/2$ for Virasoro, and $\kappa = 5c/6$ for $\cW_3$.

The Hodge class $\lambda_g := c_1(\mathbb{E}_g) \in
H^2(\Mbar_g, \mathbb{Q})$ is the first Chern class of the Hodge
bundle $\mathbb{E}_g = \pi_*\omega_{\pi}$.  The
\emph{$\hat{A}$-coefficients} $\lambda_g^{\mathrm{FP}}$ are
the genus-$g$ coefficients in the generating function
$\sum_{g \geq 1} \lambda_g^{\mathrm{FP}}\,\hbar^{2g}
= \frac{\hbar/2}{\sin(\hbar/2)} - 1$,
so that $\lambda_1^{\mathrm{FP}} = 1/24$ and
$\lambda_2^{\mathrm{FP}} = 7/5760$.


\subsection{Outline}

Section~\ref{sec:bar} reviews the bar construction for chiral
algebras and defines chirally Koszul vertex algebras.
Section~\ref{sec:mc} introduces the modular convolution algebra
and the Maurer--Cartan element $\Theta_A$.
Section~\ref{sec:algebraicity} states and proves the algebraicity
theorem (Theorem~\ref{thm:riccati}).
Section~\ref{sec:classification} establishes the four-class
partition (Theorem~\ref{thm:dichotomy}).
Section~\ref{sec:landscape} computes explicit shadow obstruction towers for
all standard families.
Section~\ref{sec:genus2} derives the $\cW_3$ cross-channel
correction to the genus-$2$ free energy by a graph-by-graph
stable graph sum.
Section~\ref{sec:connection} treats the shadow connection,
Koszul duality, and complementarity.
Section~\ref{sec:outlook} discusses open directions.

%%% ================================================================
%%%  2. CHIRAL KOSZUL PAIRS AND THE BAR CONSTRUCTION
%%% ================================================================

\section{Chiral Koszul pairs and the bar construction}
\label{sec:bar}

\subsection{The chiral bar complex}

Let $A$ be a vertex algebra on a smooth projective curve~$X$,
equipped with a conformal vector and a positive-energy grading
$A = \bigoplus_{h \geq 0} A_h$ with $\dim A_h < \infty$.
The \emph{Ran space} $\Ran(X) = \coprod_{n \geq 1}
X^n / \Sigma_n$ parameterizes nonempty finite subsets of~$X$.
For each $n \geq 1$, the Fulton--MacPherson compactification
$\FM_n(X)$ resolves the diagonals in~$X^n$ by a sequence of
iterated real blowups along all partial diagonals, producing a
smooth manifold-with-corners whose boundary strata are indexed
by trees of collisions~\cite{GK}.

\begin{definition}[Chiral bar complex]
\label{def:bar-complex}
The \emph{chiral bar complex} $B(A)$ is the factorization
coalgebra on $\Ran(X)$ with:
\begin{enumerate}[(i)]
\item \emph{Underlying space}: the graded vector space
  $B(A)_n := (s^{-1}\overline{A})^{\otimes n}$, where
  $\overline{A} = A / k\mathbf{1}$ is the augmentation ideal
  and $s^{-1}$ denotes desuspension (degree~$-1$).

\item \emph{Bar differential}: the coderivation
  $d_B\colon B(A)_n \to B(A)_n$ whose components are the
  collision residues
  \[
    d_B|_{B(A)_n \to B(A)_{n-1}}
    \;=\; \sum_{i < j} \Res_{z_i \to z_j}\,
    a_i(z_i)\,a_j(z_j)\,d\!\log(z_i - z_j),
  \]
  extracting the residue of the OPE along the logarithmic
  differential $d\!\log(z_i - z_j)$.

\item \emph{Coproduct}: the cofactorization structure
  $\Delta\colon B(A)_n \to \bigoplus_{S \sqcup T = [n]}
  B(A)_{|S|} \otimes B(A)_{|T|}$
  induced by the partition of points into clusters.
\end{enumerate}
The identity $d_B^2 = 0$ encodes the associativity of the OPE.
\end{definition}

\begin{remark}[The $d\!\log$ measure]
\label{rem:dlog-measure}
The bar differential extracts residues along $d\!\log(z-w)$,
not along $dz/(z-w)$.  The logarithmic differential absorbs
one power of $(z-w)$: if the OPE has poles at
$(z-w)^{-n-1}$, the collision residue has poles at
$(z-w)^{-n}$.  For Virasoro, the OPE has poles at orders
$4$, $2$, $1$; the bar $r$-matrix has poles at orders
$3$, $1$ only.
\end{remark}

\subsection{Chirally Koszul algebras}

\begin{definition}[Chirally Koszul]
\label{def:chirally-koszul}
A vertex algebra~$A$ is \emph{chirally Koszul} if the bar
cohomology $H^{\bullet}(B(A), d_B)$ is concentrated in bar
degree~$1$: for every $n \geq 2$, the map
$H^{\bullet}(B(A)_n) \to H^{\bullet}(B(A)_1)^{\otimes n}$
induced by the coproduct is a quasi-isomorphism.
\end{definition}

Chirally Koszul algebras are the vertex-algebraic analogues
of Koszul algebras in the sense of Priddy: the bar resolution
has a minimal model.  The standard families (Heisenberg,
affine Kac--Moody at non-critical level, $\beta\gamma$,
Virasoro, and $\cW_N$) are all chirally Koszul
(the universal affine vertex algebra $V_k(\fg)$ is Koszul for
all $k \neq -h^\vee$ by PBW concentration; the Virasoro
and $\cW_N$ algebras are Koszul by the same argument
applied to their free-field realizations).

\section{The modular convolution algebra and the Maurer--Cartan element}
\label{sec:mc}

\subsection{The modular cyclic deformation complex}

Let $A$ be a vertex algebra on a smooth projective curve~$X$
as in Section~\ref{sec:bar}.  The bar complex $B(A)$ carries,
for each stable pair $(g,n)$ with $2g - 2 + n > 0$, a
genus-$g$, $n$-marked component $B^{(g,n)}(A)$.

\begin{definition}[Modular cyclic deformation complex]
\label{def:mod-cyc}
The \emph{modular cyclic deformation complex} of~$A$ is
\begin{equation}\label{eq:mod-cyc}
  \Defcyc^{\mathrm{mod}}(A)
  \;:=\;
  \prod_{\substack{g,n \\ 2g-2+n>0}}
  \operatorname{CoDer}^{\mathrm{cyc}}\!\bigl(
    B^{(g,n)}(A)
  \bigr)[1],
\end{equation}
where the product ranges over all stable pairs.
The differential is $d = d_{\mathrm{int}} + d_{\mathrm{sew}}$:
\begin{enumerate}[(i)]
\item $d_{\mathrm{int}}$ is the internal cyclic coderivation
  differential on each component;
\item $d_{\mathrm{sew}}$ is the clutching operation induced by
  the boundary maps
  $\partial\colon \Mbar_{g,n} \supset
  \Delta_{g_1,S_1;\,g_2,S_2} \to
  \Mbar_{g_1,|S_1|+1} \times
  \Mbar_{g_2,|S_2|+1}$.
\end{enumerate}
The Lie bracket is \emph{graph composition}: for coderivations
$\xi$ and $\eta$ at genera $g_1$, $g_2$ and arities $n_1$, $n_2$,
\begin{equation}\label{eq:graph-bracket}
  [\xi,\eta]
  \;:=\;
  \sum_{\Gamma}\xi \circ_\Gamma \eta,
\end{equation}
summed over stable graphs~$\Gamma$ connecting an output of~$\xi$
to an input of~$\eta$ with cross-polarization on internal edges.
\end{definition}

The genus grading induces a complete descending filtration
$F^{\bullet}\Defcyc^{\mathrm{mod}}(A)$ with $F^g$
consisting of elements supported at genus $\geq g$.
At $g = 0$ the complex reduces to the ordinary cyclic
deformation complex of~$A$.  The bracket respects the
arity filtration: $[\xi, \eta]$ has arity $\leq n_1 + n_2 - 2$
(one input and one output are consumed by the internal edge).


\subsection{The Maurer--Cartan element}

The bar differential determines a canonical Maurer--Cartan element.

\begin{proposition}[Bar-intrinsic MC element]
\label{prop:bar-mc}
Let $D_A$ denote the full bar differential on $B(A)$
(incorporating contributions from all genera and all boundary
strata of~$\Mbar_{g,n}$), and let $d_0$ denote its genus-$0$
linear part.  Define
\begin{equation}\label{eq:theta-def}
  \Theta_A \;:=\; D_A - d_0
  \;\in\; \Defcyc^{\mathrm{mod}}(A).
\end{equation}
Then $\Theta_A$ satisfies the Maurer--Cartan equation
\begin{equation}\label{eq:mc-equation}
  d_0\,\Theta_A + \tfrac{1}{2}[\Theta_A, \Theta_A] = 0.
\end{equation}
\end{proposition}

\begin{proof}
The identity $D_A^2 = 0$ is a formal consequence of
$\partial^2 = 0$ on~$\Mbar_{g,n}$: the square of the bar
differential decomposes into a sum over pairs of boundary
strata whose contributions cancel by the codimension-$2$
face relation.  Substituting $D_A = d_0 + \Theta_A$ into
$D_A^2 = 0$ and expanding:
\[
  0 \;=\; d_0^2 + d_0\,\Theta_A + \Theta_A \circ d_0
  + \Theta_A \circ \Theta_A.
\]
Since $d_0^2 = 0$ (the genus-$0$ bar differential squares to
zero by OPE associativity) and
$\Theta_A \circ d_0 + \Theta_A \circ \Theta_A
= d_0\,\Theta_A + \tfrac{1}{2}[\Theta_A, \Theta_A]$
in the cyclic coderivation Lie algebra, the result follows.
\end{proof}


\subsection{The shadow obstruction tower}

The MC element $\Theta_A$ encodes all genera and arities.
Its finite-order truncations form the shadow obstruction tower.

\begin{definition}[Shadow obstruction tower]
\label{def:shadow-tower}
The \emph{shadow obstruction tower} of~$A$ is the sequence of truncations
$\Theta_A^{\leq r}$ ($r \geq 2$) obtained by retaining only
the components of arity at most~$r$ in $\Defcyc^{\mathrm{mod}}(A)$.
The \emph{arity-$r$ shadow} is the projection
\[
  \Sh_r(A) \;:=\;
  \pi_r(\Theta_A)
  \;\in\;
  \operatorname{CoDer}^{\mathrm{cyc}}(B^{(\bullet,r)}(A))[1],
\]
the component of $\Theta_A$ at arity exactly~$r$.
\end{definition}

The MC equation~\eqref{eq:mc-equation}, projected onto
the arity-$r$ component, yields the \emph{all-arity master
equation}:

\begin{proposition}[All-arity master equation]
\label{prop:master-eq}
For each $r \geq 2$, the arity-$r$ projection of the MC
equation gives
\begin{equation}\label{eq:master-eq}
  d_0(\Sh_r) + \sum_{\substack{j+k = r+2 \\ j,k \geq 2}}
  [\Sh_j, \Sh_k] \;=\; 0.
\end{equation}
In particular, $\Sh_r$ is determined recursively by
$\Sh_2, \ldots, \Sh_{r-1}$ together with the
vanishing of a cohomological obstruction.
\end{proposition}

\begin{proof}
The differential $d_0$ and the bracket $[-,-]$ on
$\Defcyc^{\mathrm{mod}}(A)$ both respect the arity filtration.
Restricting the MC equation to arity~$r$ gives exactly
the stated identity: the linear term is $d_0(\Sh_r)$, and
the quadratic term $\tfrac{1}{2}[\Theta_A, \Theta_A]$ at
arity~$r$ receives contributions from pairs $(j,k)$ with
$j + k = r + 2$ (the bracket consumes two inputs via the
internal edge, reducing total arity by~$2$).
\end{proof}

\begin{definition}[Primary line and shadow data]
\label{def:primary-line}
A \emph{primary line} in $\Defcyc^{\mathrm{mod}}(A)$ is a
one-dimensional subspace $L = k \cdot \eta$ spanned by a
single primary field~$\eta$ of definite conformal weight.
On~$L$, the arity-$r$ shadow restricts to a scalar:
$\Sh_r|_L = S_r \cdot x^r$, where $x$ is the coordinate
on~$L$ and $S_r = S_r(A) \in \mathbb{Q}(c)$.

The first three scalars are:
\begin{itemize}
\item $S_2 =: \kappa$, the \emph{modular characteristic}:
  $F_1(A) = \kappa/24$.
\item $S_3 =: \alpha$, the \emph{cubic shadow}:
  the three-point collision residue on $\FM_3(X)$.
\item $S_4$, the \emph{quartic shadow}:
  the four-point residue encoding the first nonlinear OPE
  invariant.
\end{itemize}
\end{definition}


\subsection{The shadow metric}

\begin{definition}[Shadow metric and critical discriminant]
\label{def:shadow-metric}
Let $L$ be a primary line with shadow data $S_r$, curvature
$\kappa = S_2 \neq 0$, and cubic coefficient $\alpha = S_3$.
Define the \emph{weighted initial data}
$a_0 := 2\kappa$, $a_1 := 3\alpha$, $a_2 := 4S_4$.
The \emph{shadow metric} on~$L$ is the quadratic polynomial
\begin{equation}\label{eq:shadow-metric}
  Q(t)
  \;:=\;
  a_0^2 + 2a_0 a_1\,t + (a_1^2 + 2a_0 a_2)\,t^2
  \;=\;
  4\kappa^2 + 12\kappa\alpha\,t
  + (9\alpha^2 + 16\kappa S_4)\,t^2.
\end{equation}
The \emph{critical discriminant} of~$L$ is
\begin{equation}\label{eq:critical-disc}
  \Delta \;:=\; 8\kappa S_4 \;=\; a_0 a_2.
\end{equation}
\end{definition}

\begin{lemma}[Gaussian decomposition]
\label{lem:gaussian}
\begin{equation}\label{eq:gaussian-decomp}
  Q(t)
  \;=\;
  (2\kappa + 3\alpha\,t)^2
  \;+\;
  2\Delta\,t^2.
\end{equation}
The first term $(a_0 + a_1 t)^2$ is the shadow metric of
the truncated tower with $S_r = 0$ for $r \geq 4$.
The second term $2\Delta\,t^2$ is the interaction correction.
\end{lemma}

\begin{proof}
$(2\kappa + 3\alpha t)^2 = 4\kappa^2 + 12\kappa\alpha\,t
+ 9\alpha^2\,t^2$.
Adding $2\Delta\,t^2 = 16\kappa S_4\,t^2$ recovers $Q(t)$.
\end{proof}


%%% ================================================================
%%%  4. PROOF OF ALGEBRAICITY
%%% ================================================================

\section{The Riccati algebraicity theorem}
\label{sec:algebraicity}

\subsection{Statement}

\begin{theorem}[Riccati algebraicity]
\label{thm:riccati}
Let $A$ be a chirally Koszul vertex algebra, $L$ a primary line
in $\Defcyc^{\mathrm{mod}}(A)$ with shadow data $S_r$, curvature
$\kappa = S_2 \neq 0$, and cubic coefficient $\alpha = S_3$.
Define the shadow metric $Q(t)$ and the weighted shadow
generating function $H(t) := \sum_{r \geq 2} r\,S_r\,t^r$.
Then
\begin{equation}\label{eq:algebraicity}
  H(t)^2 \;=\; t^4\,Q(t).
\end{equation}
Equivalently,
\begin{equation}\label{eq:closed-form}
  H(t)
  \;=\;
  t^2\sqrt{Q(t)}.
\end{equation}
Every shadow coefficient $S_r$ with $r \geq 5$ is determined by
the three genus-$0$ invariants $(\kappa, \alpha, S_4)$ via the
closed form $S_r = \frac{1}{r}\,[t^{r-2}]\sqrt{Q(t)}$.
\end{theorem}


\subsection{The single-line recursion}

The Maurer--Cartan equation on~$L$ reduces to a quadratic
convolution identity.

\begin{lemma}[MC recursion on a primary line]
\label{lem:mc-recursion}
Let $P := 1/\kappa$.  On the primary line~$L$, the master
equation \textup{(}Proposition~\textup{\ref{prop:master-eq})}
reduces to the recursion
\begin{equation}\label{eq:recursion}
  S_r
  \;=\;
  -\frac{P}{2r}
  \sum_{\substack{j + k = r + 2 \\ 3 \leq j \leq k}}
  c_{jk}\,jk\,S_j\,S_k,
  \qquad
  c_{jk} = \begin{cases} 1 & j < k, \\ 1/2 & j = k,\end{cases}
\end{equation}
for all $r \geq 5$.
\end{lemma}

\begin{proof}
On the one-dimensional subspace~$L$, the bracket
$[\Sh_j, \Sh_k]$ restricts to a scalar pairing.
The cyclic coderivation bracket on~$L$ evaluates as
$[\Sh_j, \Sh_k]_L = c_{jk}\,jk\,P\,S_j S_k\,x^{j+k-2}$:
the factor $jk$ comes from the number of ways to select
one input from each coderivation, the factor~$P = 1/\kappa$
comes from the propagator on the internal edge
(the inverse of the curvature pairing), and $c_{jk}$
corrects for the symmetry when $j = k$.
The linear term $d_0(\Sh_r)$ contributes $2r\,\kappa\,S_r$
(the eigenvalue of the curvature operator on the arity-$r$
component).  Dividing through by $2r\,\kappa$ gives the
stated recursion.

Since $j \geq 3$ forces $k = r + 2 - j \leq r - 1$,
the right-hand side involves only $S_3, \ldots, S_{r-1}$,
and the recursion is well-defined.
\end{proof}


\subsection{Proof of algebraicity}

\begin{proof}[Proof of Theorem~\textup{\ref{thm:riccati}}]
Define $a_n := (n+2)\,S_{n+2}$ for $n \geq 0$, so that
$a_0 = 2\kappa$, $a_1 = 3\alpha$, $a_2 = 4S_4$, and
$H(t) = \sum_{n \geq 0} a_n\,t^{n+2}$.
Set $F(t) := H(t)/t^2 = \sum_{n \geq 0} a_n\,t^n$.
The claim~\eqref{eq:algebraicity} is equivalent to
$F(t)^2 = Q(t)$.

Expand $F(t)^2 = \sum_{m \geq 0}
\bigl(\sum_{i+j=m} a_i a_j\bigr)\,t^m$.
The first three coefficients are:
\begin{alignat*}{2}
  m &= 0: &\quad a_0^2 &= 4\kappa^2, \\
  m &= 1: &\quad 2a_0 a_1 &= 12\kappa\alpha, \\
  m &= 2: &\quad a_1^2 + 2a_0 a_2 &= 9\alpha^2 + 16\kappa S_4.
\end{alignat*}
These are exactly the coefficients of $Q(t)$
(Definition~\ref{def:shadow-metric}).

For $m \geq 3$, we must show that the convolution
$\sum_{i+j=m} a_i a_j = 0$.
Setting $r = m + 2$, the convolution reads
\[
  \sum_{i+j=m} a_i a_j
  \;=\;
  \sum_{\substack{i+j = r-2 \\ i, j \geq 0}}
  (i+2)(j+2)\,S_{i+2}\,S_{j+2}
  \;=\;
  \sum_{\substack{p+q = r+2 \\ p, q \geq 2}}
  pq\,S_p\,S_q,
\]
where $p = i+2$, $q = j+2$.
Splitting off the terms involving $S_2 = \kappa$:
the sum over $p + q = r + 2$ with $p, q \geq 2$ decomposes
as the $p = 2$ term (contributing $2r\,\kappa\,S_r$),
the $q = 2$ term (contributing $2r\,S_r\,\kappa$ again),
and the sum with $p, q \geq 3$.  So:
\begin{equation}\label{eq:convolution-split}
  \sum_{i+j=m} a_i a_j
  \;=\;
  4r\,\kappa\,S_r
  \;+\;
  \sum_{\substack{p+q=r+2 \\ p,q \geq 3}} pq\,S_p\,S_q.
\end{equation}
The MC recursion (Lemma~\ref{lem:mc-recursion}) gives
\[
  S_r = -\frac{1}{2r\kappa}
  \sum_{\substack{p+q=r+2 \\ 3 \leq p \leq q}}
  c_{pq}\,pq\,S_p\,S_q.
\]
Writing the symmetrized sum:
$\sum_{\substack{p+q=r+2 \\ p,q \geq 3}} pq\,S_p\,S_q
= 2 \sum_{\substack{p+q=r+2 \\ 3 \leq p < q}} pq\,S_p S_q
+ \sum_{\substack{p+q=r+2 \\ p=q}} p^2\,S_p^2
= 2 \sum_{\substack{p+q=r+2 \\ 3 \leq p \leq q}}
c_{pq}\,pq\,S_p S_q
= -4r\,\kappa\,S_r$.
Substituting into~\eqref{eq:convolution-split}:
\[
  \sum_{i+j=m} a_i a_j
  \;=\;
  4r\,\kappa\,S_r + (-4r\,\kappa\,S_r)
  \;=\; 0.
\]
Hence $[t^m](F^2 - Q) = 0$ for all $m \geq 3$.
Combined with the verification at $m = 0, 1, 2$,
this proves $F(t)^2 = Q(t)$, hence $H(t)^2 = t^4\,Q(t)$.

The closed form for $S_r$ follows:
$r\,S_r = [t^{r-2}]F(t) = [t^{r-2}]\sqrt{Q(t)}$.
\end{proof}

\begin{remark}[Riccati ODE reformulation]
\label{rem:riccati-ode}
Set $U(t) := \sum_{r \geq 3} S_r\,t^r$, the nonlinear part
of the shadow generating function (removing $\Sh_2$).
The algebraic relation~\eqref{eq:algebraicity} is equivalent
to the first-order ODE
\begin{equation}\label{eq:riccati-ode}
  2t\,U'(t) \;+\; \frac{1}{2\kappa}\bigl(U'(t)\bigr)^2
  \;=\; 6\alpha\,t^3
  + \frac{2(\alpha^2 + 2\Delta)}{\kappa}\,t^4,
\end{equation}
which is \emph{quadratic in~$U'$}.  This is the precise sense
in which the shadow obstruction tower is algebraic of degree~$2$: the
generating function satisfies a Riccati-type equation,
and all terms at $t^r$ for $r \geq 5$ cancel identically.
The ODE determines $U$ uniquely from $(\kappa, \alpha, S_4)$
by power-series inversion.
\end{remark}


\subsection{The shadow connection}

\begin{definition}[Shadow connection]
\label{def:shadow-conn}
The \emph{shadow connection} on $L \setminus \{Q = 0\}$ is
\begin{equation}\label{eq:shadow-conn}
  \nabla^{\mathrm{sh}}
  \;:=\; d \;-\; \frac{Q'(t)}{2\,Q(t)}\,dt
  \;=\; d - d\!\left(\log\sqrt{Q}\right).
\end{equation}
\end{definition}

\begin{proposition}[Properties of the shadow connection]
\label{prop:conn-properties}
\leavevmode
\begin{enumerate}[(i)]
\item \emph{Flatness.}\enspace
  $(\nabla^{\mathrm{sh}})^2 = 0$, since
  $Q'/(2Q)\,dt$ is closed on a one-dimensional base.

\item \emph{Regular singularities.}\enspace
  At each zero $t_0$ of~$Q$, the connection form has
  a simple pole with residue~$\tfrac{1}{2}$:
  near~$t_0$, $Q(t) \approx Q'(t_0)(t - t_0)$, so
  $\frac{Q'}{2Q}\,dt \approx \frac{1}{2(t-t_0)}\,dt$.

\item \emph{Koszul monodromy.}\enspace
  The monodromy around each branch point is
  $\exp(2\pi i \cdot \tfrac{1}{2}) = -1$.
  The monodromy representation factors through
  $\mathbb{Z}/2$; the nontrivial element acts by~$-1$.

\item \emph{Flat section.}\enspace
  The unique flat section with $\Phi(0) = 1$ is
  \begin{equation}\label{eq:flat-section}
    \Phi(t)
    \;=\; \sqrt{Q(t)/Q(0)}
    \;=\; \frac{\sqrt{Q(t)}}{2\kappa}.
  \end{equation}
  The shadow generating function is the flat section
  weighted by the arity offset:
  \begin{equation}\label{eq:H-from-Phi}
    H(t) \;=\; 2\kappa\,t^2\,\Phi(t).
  \end{equation}
\end{enumerate}
\end{proposition}

\begin{proof}
Parts (i)--(iii) are standard properties of the Gauss--Manin
connection of the family $\{\sqrt{Q(t)}\}_{t \in L}$:
the unique flat connection for which $\sqrt{Q}$ is a horizontal
section.  The residue computation at a simple zero of~$Q$ is
immediate from the Laurent expansion.  The monodromy is that
of a square root around a simple branch point.

Part~(iv): $\nabla^{\mathrm{sh}}\Phi
= d\Phi - \frac{Q'}{2Q}\Phi\,dt = 0$ is verified directly,
since $d\Phi = \frac{Q'}{2\sqrt{Q \cdot Q(0)}}\,dt
= \frac{Q'}{2Q}\,\Phi\,dt$.
The identity $H = 2\kappa\,t^2\,\Phi$ is
Theorem~\ref{thm:riccati}: $H = t^2\sqrt{Q}
= t^2 \cdot 2\kappa \cdot \Phi$.
\end{proof}

\begin{remark}[Parallel transport of curvature]
\label{rem:parallel-transport}
At $t = 0$, the shadow reduces to the curvature~$\kappa$.
The flat section $\Phi(t) = \sqrt{Q(t)}/(2\kappa)$ incorporates
higher OPE data through~$\nabla^{\mathrm{sh}}$.
The generating function $H(t) = 2\kappa\,t^2\,\Phi(t)$
is \emph{not} a flat section of $\nabla^{\mathrm{sh}}$: it has
a double zero at $t = 0$ from the arity offset~$t^2$.
\end{remark}

\begin{remark}[The Koszul monodromy]
\label{rem:koszul-monodromy}
The shadow connection is the Gauss--Manin connection of a
family of double covers: the spectral curve
$\Sigma = \{H^2 = t^4 Q(t)\}$ is a double cover of the
$t$-line, ramified at the zeros of~$Q$.
Analytic continuation of $\sqrt{Q}$ around a branch point
returns~$-\sqrt{Q}$; this is the standard $\mathbb{Z}/2$
monodromy of a square root.

Under Koszul duality $A \mapsto A^!$, the shadow data
transforms: for the Virasoro family,
$\kappa(c) \mapsto \kappa(26-c)$,
$\alpha(c) \mapsto \alpha(26-c)$,
$S_4(c) \mapsto S_4(26-c)$.  The shadow metric
$Q$ transforms by $c \mapsto 26 - c$, and the
connection transforms by the Verdier involution.
At $c = 13$ (the self-dual point):
$Q = Q^!$ and the connection is self-dual.
\end{remark}

\section{The four-class classification}
\label{sec:classification}

The algebraicity theorem (Theorem~\ref{thm:riccati}) implies
that the shadow obstruction tower is determined by three invariants
$(\kappa, \alpha, S_4)$, and that the generating function
$H(t) = t^2\sqrt{Q(t)}$ is algebraic of degree~$2$ over
the parameter line.  The depth of the tower, defined as
$r_{\max} := \sup\{r : S_r \neq 0\}$, is controlled entirely
by the factorization of the shadow metric~$Q$.

\begin{theorem}[Single-line dichotomy]
\label{thm:dichotomy}
On any primary line $L$ with $\kappa \neq 0$ and critical
discriminant $\Delta = 8\kappa S_4$, the shadow depth satisfies
$r_{\max} \in \{2, 3, \infty\}$, classified as follows.
\begin{enumerate}[(i)]
\item $\Delta = 0$, $\alpha = 0$:
  $Q = (2\kappa)^2$; $H(t) = 2\kappa\,t^2$;
  $r_{\max} = 2$ \textup{(class~$\mathbf{G}$, Gaussian)}.
\item $\Delta = 0$, $\alpha \neq 0$:
  $Q = (2\kappa + 3\alpha\,t)^2$;
  $H(t) = t^2(2\kappa + 3\alpha\,t)$;
  $r_{\max} = 3$ \textup{(class~$\mathbf{L}$, Lie)}.
\item $\Delta \neq 0$:
  $Q$ is irreducible over $\mathbb{Q}(\kappa, \alpha, S_4)$;
  $S_r \neq 0$ for all $r \geq 4$
  generically; $r_{\max} = \infty$ \textup{(class~$\mathbf{M}$, mixed)}.
\end{enumerate}
For $r \geq 4$: $S_r = \Delta \cdot R_r$ with
$R_r \in \mathbb{Q}(\alpha, \kappa, S_4)$ independent of the
algebra family.
\end{theorem}

\begin{proof}
The Gaussian decomposition
$Q(t) = (2\kappa + 3\alpha\,t)^2 + 2\Delta\,t^2$
shows that $Q$ is a perfect square if and only if
$\Delta = 0$.

\emph{Case $\Delta = 0$.}
Then $Q = (2\kappa + 3\alpha\,t)^2$, so
$H(t) = t^2(2\kappa + 3\alpha\,t)$.
If also $\alpha = 0$, then $H(t) = 2\kappa\,t^2$,
i.e., $S_r = 0$ for all $r \geq 3$.  This is class~$\mathbf{G}$.
If $\alpha \neq 0$, then $H$ is a cubic polynomial in~$t$,
so $S_r = 0$ for $r \geq 4$ while $S_3 = \alpha \neq 0$.
This is class~$\mathbf{L}$.

\emph{Case $\Delta \neq 0$.}
Then $Q$ has two distinct zeros
$t_{\pm} = \bigl(-6\kappa\alpha \pm
\sqrt{-8\kappa^2\Delta}\bigr)\big/
(9\alpha^2 + 2\Delta)$,
so $\sqrt{Q(t)}$ is not polynomial.  The Taylor coefficients
$S_r = \frac{1}{r}[t^{r-2}]\sqrt{Q(t)}$ are nonzero for all
$r \geq 4$ generically (the only way $[t^n]\sqrt{Q}$ can
vanish identically is if $\sqrt{Q}$ is polynomial, which
requires $\Delta = 0$).  Hence $r_{\max} = \infty$:
class~$\mathbf{M}$.

To see that $S_r$ is divisible by~$\Delta$ for $r \geq 4$,
write $\sqrt{Q} = |2\kappa + 3\alpha\,t|\bigl(1 +
2\Delta\,t^2/(2\kappa+3\alpha\,t)^2\bigr)^{1/2}$
and expand the binomial series.  Every term beyond zeroth
order contributes a factor of~$\Delta$, and all Taylor
coefficients of $H$ at order $\geq 4$ come from these terms.
The universal rational function $R_r$ is the coefficient
of $\Delta$ in this expansion and depends only on
$\kappa$, $\alpha$, and $S_4$.
\end{proof}

A fourth class escapes the single-line dichotomy.

\begin{definition}[Class C: contact depth]\label{def:class-C}
A chirally Koszul vertex algebra is of \emph{class~$\mathbf{C}$}
(contact, $r_{\max} = 4$) if
\begin{enumerate}[(a)]
\item on every one-dimensional primary line $L$, the shadow
  metric $Q_L$ is a perfect square ($\Delta_L = 0$),
  so the single-line tower terminates at arity~$3$; but
\item in the full cyclic deformation complex
  $\Defcyc^{\mathrm{mod}}(A)$, there exists a charged stratum
  supporting a nonzero quartic contact invariant
  $Q^{\mathrm{cont}} \neq 0$ at arity~$4$.
\end{enumerate}
The quartic contact invariant lives on a two-dimensional
weight/contact slice.  Its self-bracket under the
$H$-Poisson structure exits the single-line complex by
\emph{rank-one rigidity}: on a one-dimensional target,
a quadratic self-bracket produces a scalar multiple of the
generator, but the contact direction is orthogonal to the
primary line, so the image lies outside the complex.
The quintic obstruction $o_5$ therefore vanishes, and the
tower terminates at arity~$4$.
\end{definition}

\begin{example}[Class C: the $\beta\gamma$ system]
\label{ex:betagamma-classC}
The $\beta\gamma$ system of conformal weight $\lambda$ has
$\kappa = 6\lambda^2 - 6\lambda + 1$ (equal to $c/2$ where
$c = 2(6\lambda^2 - 6\lambda + 1)$).  On the weight-changing
primary line, $\alpha = 0$ and $S_4 = 0$ (abelian OPE along
this direction), so $\Delta = 0$ and the single-line tower
terminates.  But the full two-variable deformation complex
contains a charged stratum with a nonzero quartic contact
invariant.  The quintic obstruction vanishes by rank-one
rigidity.  Hence $\beta\gamma$ is class~$\mathbf{C}$
with $r_{\max} = 4$.
\end{example}

\begin{remark}[The four classes partition the Koszul world]
\label{rem:four-classes-partition}
Every chirally Koszul vertex algebra with $\kappa \neq 0$
falls into exactly one of the four classes G, L, C, M.
All four classes are Koszul: shadow depth classifies the
complexity of the OPE within the Koszul world, not
Koszulness itself.
The four classes have the following structural
characterization:
\begin{center}
\renewcommand{\arraystretch}{1.3}
\begin{tabular}{ccccl}
\toprule
Class & $\alpha$ & $\Delta$ & $r_{\max}$ & Examples \\
\midrule
$\mathbf{G}$ & $0$ & $0$ & $2$ & Heisenberg, lattice VOAs, free fermion \\
$\mathbf{L}$ & $\neq 0$ & $0$ & $3$ & Affine Kac--Moody (all types) \\
$\mathbf{C}$ & $0^*$ & $0^*$ & $4$ & $\beta\gamma$, $bc$ ghosts \\
$\mathbf{M}$ & $\neq 0$ & $\neq 0$ & $\infty$ & Virasoro, $\cW_N$ ($N \geq 3$) \\
\bottomrule
\end{tabular}
\end{center}
\noindent The ${}^*$ for class~$\mathbf{C}$ indicates that
$\alpha = 0$ and $\Delta = 0$ hold on every one-dimensional
primary line, while the quartic contact invariant lives on a
higher-dimensional stratum.
\end{remark}


\subsection{Shadow depth and the depth decomposition}
\label{subsec:depth-decomposition}

For algebras of class~$\mathbf{M}$, the shadow obstruction tower is
infinite.  The total depth admits a decomposition into
arithmetic and algebraic components.

\begin{definition}[Shadow depth decomposition]
\label{def:depth-decomposition}
For a chirally Koszul vertex algebra~$A$ of class~$\mathbf{M}$,
define the \emph{total shadow depth}
\[
  d(A) \;=\; 1 + d_{\mathrm{arith}} + d_{\mathrm{alg}}
\]
where:
\begin{itemize}
\item $d_{\mathrm{alg}} \in \{0, 1, 2, \infty\}$ is the
  \emph{algebraic depth}, measuring the complexity of the
  OPE-generated shadow obstruction tower on the primary line.  By the
  algebraicity theorem, $d_{\mathrm{alg}} = 0$ for all
  class~$\mathbf{M}$ algebras: three invariants determine
  the full tower.
\item $d_{\mathrm{arith}} \geq 0$ is the \emph{arithmetic
  depth}, counting contributions from cusp forms.
  Shadow coefficients of weight $\leq 2r$ at arity~$r$
  lie in the subspace spanned by Eisenstein series;
  cusp-form contributions first appear at genus~$g$ where
  $S_{2g}(\mathrm{SL}_2(\mathbb{Z})) \neq 0$, i.e., at
  $g \geq 6$ (the Ramanujan $\Delta$ form at weight~$12$).
\end{itemize}
\end{definition}

\begin{remark}
The algebraic depth is determined by genus-$0$ OPE data
(Theorem~\ref{thm:riccati}).  The arithmetic depth reflects
automorphic phenomena in the genus expansion:\ cusp-form
contributions are invisible at low genus but potentially
enter at $g \geq 6$.  For all families treated in
Section~\ref{sec:landscape}, $d_{\mathrm{arith}} = 0$
through the genera computed here.
\end{remark}


\subsection{Shadow radius and convergence}
\label{subsec:shadow-radius}

For class~$\mathbf{M}$ algebras, the shadow obstruction tower is
infinite.  The algebraicity theorem determines the growth
rate of the coefficients $S_r$ exactly.

\begin{definition}[Shadow growth rate]
\label{def:shadow-radius}
Let $A$ be a chirally Koszul vertex algebra of
class~$\mathbf{M}$ with shadow data $(\kappa, \alpha, S_4)$
and $\Delta = 8\kappa S_4 \neq 0$.  The \emph{shadow growth
rate} (or shadow radius) is
\[
  \rho(A)
  \;:=\;
  \frac{\sqrt{9\alpha^2 + 2\Delta}}{2|\kappa|}.
\]
\end{definition}

\begin{proposition}[Shadow growth rate as reciprocal convergence radius]
\label{prop:rho-convergence}
The growth rate $\rho(A)$ equals the reciprocal of the
radius of convergence of the weighted generating function
$H(t) = t^2\sqrt{Q(t)}$:
\[
  \rho(A) \;=\; \frac{1}{R},
  \qquad
  R \;=\; \min_{Q(t_0)=0} |t_0|.
\]
The branch points of $H$ are the zeros of $Q$.  When
$\Delta > 0$ (the generic case for class~$\mathbf{M}$),
these form a complex conjugate pair of equal modulus, so
\[
  R
  \;=\;
  \frac{2|\kappa|}{\sqrt{9\alpha^2 + 2\Delta}}.
\]
The shadow coefficients satisfy the asymptotic formula
\[
  S_r
  \;\sim\;
  C(A)\,\rho(A)^r\,r^{-5/2}\,\cos(r\theta + \phi),
  \qquad r \to \infty,
\]
where $\theta$ and $\phi$ are determined by the argument
of the nearest branch point, and $C(A)$ is a computable
algebraic constant.
\end{proposition}

\begin{proof}
Since $H(t) = t^2\sqrt{Q(t)}$ and $Q$ is quadratic, $H$
is an algebraic function of degree~$2$ with branch points
at the zeros of~$Q$.  The radius of convergence of the
Taylor series of $\sqrt{Q}$ equals the distance from
$t = 0$ to the nearest branch point.  The zeros of $Q$
are at
$t_{\pm} = (-6\kappa\alpha \pm \sqrt{-8\kappa^2\Delta})/
(9\alpha^2 + 2\Delta)$.
For $\Delta > 0$, the discriminant
$-8\kappa^2\Delta < 0$, so $t_{\pm}$ are complex
conjugate with common modulus
$|t_{\pm}| = 2|\kappa|/\sqrt{9\alpha^2 + 2\Delta}$.
The transfer theorem of Flajolet--Sedgewick gives the
stated asymptotic formula with exponent $-5/2$
(from the square-root singularity of an algebraic
function of degree~$2$).
\end{proof}

\begin{example}[Virasoro shadow growth rate]
\label{ex:virasoro-rho}
For Virasoro at central charge~$c > 0$ with $\kappa = c/2$,
$\alpha = 2$, $\Delta = 40/(5c+22)$:
\[
  \rho(\mathrm{Vir}_c)
  \;=\;
  \frac{1}{c}\sqrt{\frac{180c + 872}{5c + 22}}.
\]
The equation $\rho(\mathrm{Vir}_{c^*}) = 1$ is equivalent
to the cubic
\begin{equation}\label{eq:critical-cubic}
  5c^3 + 22c^2 - 180c - 872 \;=\; 0,
\end{equation}
whose unique positive real root is
$c^* \approx 6.125$.
\begin{itemize}
\item For $c > c^*$: $\rho < 1$, the shadow obstruction tower converges
  (summable amplitudes).  In particular, $\rho(13) \approx
  0.467$, $\rho(26) \approx 0.232$.
\item For $c < c^*$: $\rho > 1$, the shadow obstruction tower diverges
  (exponential growth).  In particular,
  $\rho(1) \approx 6.242$, $\rho(1/2) \approx 12.532$.
\end{itemize}
All unitary minimal models ($c < 1$) have strongly
divergent shadow obstruction towers.  The string-theoretic value
$c = 26$ has a convergent tower with $\rho \approx 0.232$.
\end{example}


%%% ================================================================
%%%  6. EXPLICIT SHADOW TOWERS FOR THE STANDARD LANDSCAPE
%%% ================================================================

\section{Explicit shadow obstruction towers for the standard landscape}
\label{sec:landscape}

The shadow invariants $(\kappa, \alpha, S_4,
\Delta, \rho)$ for all standard families are as follows.

\subsection{Class G: Gaussian algebras}
\label{subsec:class-G}

\begin{proposition}[Heisenberg shadow obstruction tower]
\label{prop:heisenberg-shadow}
For the Heisenberg vertex algebra $\cH_k$ of rank~$1$
at level $k \neq 0$:
\[
  \kappa \;=\; k,
  \qquad
  \alpha \;=\; 0,
  \qquad
  S_r \;=\; 0 \quad\text{for all } r \geq 3.
\]
The shadow obstruction tower terminates immediately:
$\Theta_{\cH_k}^{\leq 2} = k\,\eta^{\otimes 2}$,
$Q(t) = 4k^2$ \textup{(}constant\textup{)}, $H(t) = 2k\,t^2$.
The genus-$g$ free energy is $F_g(\cH_k) = k \cdot
\lambda_g^{\mathrm{FP}}$ for all $g \geq 1$.
\end{proposition}

\begin{proof}
The Heisenberg OPE $J(z)J(w) \sim k/(z-w)^2$ is purely
quadratic: all higher collision residues on $\FM_n(X)$
for $n \geq 3$ vanish identically.  Hence $S_r = 0$ for
$r \geq 3$, $\alpha = 0$, $S_4 = 0$, $\Delta = 0$.
The shadow metric is the constant $Q = 4k^2$, which
is a perfect square.  Class~$\mathbf{G}$.
\end{proof}

\begin{proposition}[Lattice VOA shadow obstruction tower]
\label{prop:lattice-shadow}
For the lattice vertex algebra $V_\Lambda$ associated to
an even positive-definite lattice $\Lambda$ of rank~$r$:
\[
  \kappa \;=\; r,
  \qquad
  \alpha \;=\; 0,
  \qquad
  S_n \;=\; 0 \quad\text{for all } n \geq 3.
\]
Class~$\mathbf{G}$, independent of the lattice cocycle.
\end{proposition}

\begin{proof}
The primary line of a lattice VOA is spanned by the
Heisenberg currents $J^i$, $i = 1, \ldots, r$.
The OPE $J^i(z)J^j(w) \sim \delta^{ij}/(z-w)^2$
is abelian (no cubic or higher terms on the primary
line).  The lattice vertex operators $e^\alpha$ have
nontrivial OPEs but do not contribute to the primary-line
shadow obstruction tower: they live on charged strata.  The
modular characteristic $\kappa = r$ equals the
rank, independent of the specific lattice structure
(even self-dual, Niemeier, Leech, etc.).
\end{proof}

\begin{remark}[Free fermion]
\label{rem:free-fermion}
The free fermion $\psi$ with $\psi(z)\psi(w) \sim
1/(z-w)$ has $c = 1/2$ and $\kappa = c/2 = 1/4$.
On the single-generator primary line, the OPE is
purely quadratic (Clifford), so $\alpha = 0$, $S_4 = 0$,
$\Delta = 0$.  Class~$\mathbf{G}$, $r_{\max} = 2$.
\end{remark}

\subsection{Class L: Lie algebras}
\label{subsec:class-L}

\begin{proposition}[Affine Kac--Moody shadow obstruction tower]
\label{prop:affine-shadow}
For the universal affine vertex algebra $V_k(\fg)$ at
level $k \neq -h^\vee$ with $\fg$ simple of
dimension $\dim(\fg)$ and dual Coxeter number $h^\vee$:
\[
  \kappa
  \;=\;
  \frac{\dim(\fg)\,(k + h^\vee)}{2h^\vee},
  \qquad
  \alpha \;\neq\; 0,
  \qquad
  S_4 \;=\; 0,
  \qquad
  \Delta \;=\; 0.
\]
The cubic shadow $\alpha$ is nonzero (generated by the
Lie bracket), and the quartic shadow $S_4$ vanishes on
the primary line (killed by the Jacobi identity).
The shadow obstruction tower terminates at arity~$3$:
$r_{\max} = 3$, class~$\mathbf{L}$.
The shadow metric is the perfect square
$Q(t) = (2\kappa + 3\alpha\,t)^2$.
\end{proposition}

\begin{proof}
The quartic collision residue on $\FM_4(X)$, restricted
to the single-generator primary line, is proportional to
$\sum_{\textup{cyclic}} f^{abx}f^{xcd}$,
which vanishes by the Jacobi identity
$[[\cdot,\cdot],\cdot] + \textup{cyclic} = 0$.
Hence $S_4 = 0$, $\Delta = 8\kappa \cdot 0 = 0$, and
$Q$ is a perfect square.
The cubic shadow $\alpha$ encodes the Lie bracket
structure constants (the Killing $3$-cocycle on $\fg$)
and is nonzero on generic primary directions.
\end{proof}

The following table gives $\kappa$ for the simple Lie
algebras at level $k = 1$:

\begin{center}
\renewcommand{\arraystretch}{1.2}
\begin{tabular}{lccccc}
\toprule
$\fg$ & $\dim(\fg)$ & $h^\vee$ &
  $\kappa(k)$ & $\kappa(1)$ & $c(1)$ \\
\midrule
$\mathfrak{sl}_2$  &  $3$  & $2$ &
  $\frac{3(k+2)}{4}$        & $\frac{9}{4}$    & $1$ \\[2pt]
$\mathfrak{sl}_3$  &  $8$  & $3$ &
  $\frac{4(k+3)}{3}$        & $\frac{16}{3}$   & $2$ \\[2pt]
$\mathfrak{sl}_N$  &  $N^2{-}1$ & $N$ &
  $\frac{(N^2{-}1)(k+N)}{2N}$ & --- & --- \\[2pt]
$\mathfrak{so}_5$  & $10$ & $3$ &
  $\frac{10(k+3)}{6}$       & $\frac{20}{3}$ & $\frac{5}{2}$ \\[2pt]
$G_2$              & $14$ & $4$ &
  $\frac{14(k+4)}{8}$       & $\frac{35}{4}$ & $\frac{14}{5}$ \\[2pt]
$E_8$              & $248$ & $30$ &
  $\frac{248(k+30)}{60}$ & $\frac{1922}{15}$ & $8$ \\[2pt]
\bottomrule
\end{tabular}
\end{center}

\noindent For all types and all non-critical levels
$k \neq -h^\vee$: $\alpha \neq 0$, $S_4 = 0$,
$\Delta = 0$, class~$\mathbf{L}$.
The Koszul dual $V_k(\fg)^{!} = \mathrm{CE}^{\mathrm{ch}}(\fg_{-k-2h^\vee})$
has modular characteristic $\kappa(V_k(\fg)^!) = -\kappa(V_k(\fg))$,
so $\kappa(k) + \kappa(k') = 0$ (exact anti-symmetry).  Note:
$V_k(\fg)^!$ is the chiral CE algebra, not $V_{-k-2h^\vee}(\fg)$ itself.


\subsection{Class C: contact algebras}
\label{subsec:class-C}

\begin{proposition}[$\beta\gamma$ shadow obstruction tower]
\label{prop:betagamma-shadow}
For the $\beta\gamma$ system of conformal weight $\lambda$:
\begin{alignat*}{2}
  &\kappa \;=\; 6\lambda^2 - 6\lambda + 1
  \;=\; \tfrac{c}{2},
  \qquad
  &&c \;=\; 2(6\lambda^2 - 6\lambda + 1), \\
  &\alpha \;=\; 0
  &&\textup{(weight-changing line)}, \\
  &S_4 \;=\; 0
  &&\textup{(weight-changing line)}, \\
  &\Delta \;=\; 0
  &&\textup{(weight-changing line)}.
\end{alignat*}
On the weight-changing primary line, the tower terminates
at arity~$2$ (class~$\mathbf{G}$ restricted to this line).
The overall classification is class~$\mathbf{C}$: a
nonzero quartic contact invariant lives on the charged
stratum of the full two-variable deformation complex,
and the quintic vanishes by rank-one rigidity.
\end{proposition}

\begin{remark}[$bc$ ghosts]
\label{rem:bc-ghosts}
The $bc$ ghost system of spin $j$ is the fermionic mirror
of $\beta\gamma$, with $\kappa_{bc} = -(6j^2 - 6j + 1)$.
At $j = 2$ (the reparametrization ghosts):
$\kappa = -13$, $c = -26$.  Same class~$\mathbf{C}$.
\end{remark}


\subsection{Class M: mixed algebras}
\label{subsec:class-M}

\begin{proposition}[Virasoro shadow obstruction tower]
\label{prop:virasoro-shadow}
For the Virasoro vertex algebra at central charge $c \neq 0$,
$c \neq -22/5$:
\[
  \kappa \;=\; \frac{c}{2},
  \qquad
  \alpha \;=\; 2,
  \qquad
  S_4 \;=\; \frac{10}{c(5c+22)},
  \qquad
  \Delta \;=\; \frac{40}{5c+22}.
\]
The shadow metric is
\[
  Q(t)
  \;=\;
  (c + 6t)^2 + \frac{80\,t^2}{5c + 22},
\]
the tower is infinite \textup{(}class~$\mathbf{M}$\textup{)},
and the shadow growth rate is
$\rho = \frac{1}{c}\sqrt{(180c + 872)/(5c+22)}$.
\end{proposition}

\begin{theorem}[Virasoro shadow coefficients]
\label{thm:virasoro-coefficients}
The Virasoro shadow obstruction tower coefficients through arity~$10$
are:
\begin{align*}
  S_2 &\;=\; \frac{c}{2},
  \\[3pt]
  S_3 &\;=\; 2,
  \\[3pt]
  S_4 &\;=\; \frac{10}{c(5c+22)},
  \\[3pt]
  S_5 &\;=\; \frac{-48}{c^2(5c+22)},
  \\[3pt]
  S_6 &\;=\; \frac{80(45c+193)}{3\,c^3(5c+22)^2},
  \\[3pt]
  S_7 &\;=\; \frac{-2880(15c+61)}{7\,c^4(5c+22)^2},
  \\[3pt]
  S_8 &\;=\; \frac{80(2025c^2+16470c+33314)}{c^5(5c+22)^3},
  \\[3pt]
  S_9 &\;=\; \frac{-1280(2025c^2+15570c+29554)}{3\,c^6(5c+22)^3},
  \\[3pt]
  S_{10} &\;=\; \frac{256(91125c^3+1050975c^2+3989790c+4969967)}{c^7(5c+22)^4}.
\end{align*}
All coefficients belong to $\mathbb{Q}(c)$ with poles
only at $c = 0$ and $c = -22/5$.  The signs alternate
starting from $S_5$:
$S_5 < 0$, $S_6 > 0$, $S_7 < 0$, \ldots\ for $c > 0$.
\end{theorem}

\begin{proof}
The computation proceeds inductively using the master
equation.  The Hessian propagator on the primary line
is $P = 2/c$ (inverse of $\kappa = c/2$).
The arity-$r$ obstruction is
\[
  o_r
  \;=\; \sum_{\substack{j + k = r + 2 \\ j, k \geq 2}}
  \epsilon_{jk}\,\{S_j\,x^j,\, S_k\,x^k\}_H,
  \qquad
  \{f,g\}_H \;=\; f'(x)\,P\,g'(x),
\]
where $\epsilon_{jk} = 1$ for $j \neq k$ and
$\epsilon_{jk} = 1/2$ for $j = k$.
The shadow coefficient is $S_r = -o_r/(2r)$, where we
extract the coefficient of $x^r$.

Each coefficient agrees with the closed-form extraction
$S_r = \frac{1}{r}[t^{r-2}]\sqrt{Q_{\mathrm{Vir}}(t)}$
from Theorem~\ref{thm:riccati}.
\end{proof}

Evaluations at three central charges:

\begin{center}
\renewcommand{\arraystretch}{1.2}
\begin{tabular}{rccc}
\toprule
$r$ & $S_r(c{=}1)$ & $S_r(c{=}13)$ & $S_r(c{=}26)$ \\
\midrule
$2$ & $1/2$ & $13/2$ & $13$ \\
$3$ & $2$ & $2$ & $2$ \\
$4$ & $10/27$ & $10/1131$ & $5/1976$ \\
$5$ & $-16/9$ & $-16/4901$ & $-3/6422$ \\
$6$ & $19040/2187$ & $62240/49887279$ & $6815/76139232$ \\
$7$ & $-24320/567$ & $-81920/168138607$ & $-20295/1154778352$ \\
$8$ & $4144720/19683$ & $47171920/244497554379$
  & $4576085/1303909727744$ \\
\bottomrule
\end{tabular}
\end{center}

\noindent At $c = 26$ the coefficients decay rapidly
($\rho \approx 0.232$); at $c = 13$ moderately
($\rho \approx 0.467$); at $c = 1$ they grow
exponentially ($\rho \approx 6.242$).


\begin{proposition}[$\cW_N$ shadow data]
\label{prop:wN-shadow}
For the $\cW_N$ algebra ($N \geq 3$) at central charge~$c$,
the modular characteristic on the T-line (Virasoro direction)
is $\kappa_T = c/2$ with shadow data identical to
Virasoro.  The total modular characteristic is
\[
  \kappa(\cW_N)
  \;=\; \rho_N \cdot c,
  \qquad
  \rho_N \;:=\; H_N - 1
  \;=\; \sum_{j=2}^{N} \frac{1}{j},
\]
where $H_N = \sum_{j=1}^{N} 1/j$ is the $N$-th harmonic
number.
\end{proposition}

The anomaly ratios $\rho_N$ for small~$N$ are:

\begin{center}
\renewcommand{\arraystretch}{1.2}
\begin{tabular}{cccc}
\toprule
$N$ & $\cW_N$ & $\rho_N = H_N - 1$ & $\kappa(\cW_N)$ \\
\midrule
$2$ & Virasoro & $1/2$ & $c/2$ \\
$3$ & $\cW_3$ & $5/6$ & $5c/6$ \\
$4$ & $\cW_4$ & $13/12$ & $13c/12$ \\
$5$ & $\cW_5$ & $77/60$ & $77c/60$ \\
$6$ & $\cW_6$ & $29/20$ & $29c/20$ \\
$7$ & $\cW_7$ & $223/140$ & $223c/140$ \\
\bottomrule
\end{tabular}
\end{center}

\noindent All $\cW_N$ for $N \geq 3$ are class~$\mathbf{M}$:
the T-line shadow data coincides with Virasoro
($\alpha_T = 2$, $\Delta_T = 40/(5c+22)$), and the W-line
carries independent class~$\mathbf{M}$ data with even-arity
shadow obstruction tower.


\subsection{Master table}
\label{subsec:master-table}

The following table collects the shadow invariants for
the entire standard landscape.

\begin{center}
\renewcommand{\arraystretch}{1.4}
\small
\begin{tabular}{lccccccc}
\toprule
Family & $\kappa$ & $\alpha$ & $S_4$ & $\Delta$
  & Class & $r_{\max}$ & $\rho$ \\
\midrule
$\cH_k$ (Heisenberg)
  & $k$ & $0$ & $0$ & $0$ & $\mathbf{G}$ & $2$ & $0$ \\[2pt]
$V_\Lambda$ (lattice)
  & $\mathrm{rk}(\Lambda)$ & $0$ & $0$ & $0$
  & $\mathbf{G}$ & $2$ & $0$ \\[2pt]
Free fermion
  & $1/4$ & $0$ & $0$ & $0$
  & $\mathbf{G}$ & $2$ & $0$ \\[2pt]
\addlinespace
$V_k(\mathfrak{sl}_2)$
  & $\frac{3(k{+}2)}{4}$ & $\neq 0$ & $0$ & $0$
  & $\mathbf{L}$ & $3$ & $0$ \\[2pt]
$V_k(\mathfrak{sl}_N)$
  & $\frac{(N^2{-}1)(k{+}N)}{2N}$ & $\neq 0$ & $0$ & $0$
  & $\mathbf{L}$ & $3$ & $0$ \\[2pt]
$V_k(\fg)$ (general)
  & $\frac{\dim\fg\,(k{+}h^\vee)}{2h^\vee}$ & $\neq 0$ & $0$ & $0$
  & $\mathbf{L}$ & $3$ & $0$ \\[2pt]
\addlinespace
$\beta\gamma_\lambda$
  & $6\lambda^2{-}6\lambda{+}1$ & $0^*$ & $0^*$ & $0^*$
  & $\mathbf{C}$ & $4$ & --- \\[2pt]
$bc_j$
  & $-(6j^2{-}6j{+}1)$ & $0^*$ & $0^*$ & $0^*$
  & $\mathbf{C}$ & $4$ & --- \\[2pt]
\addlinespace
$\mathrm{Vir}_c$
  & $c/2$ & $2$ & $\frac{10}{c(5c{+}22)}$
  & $\frac{40}{5c{+}22}$
  & $\mathbf{M}$ & $\infty$
  & $\frac{1}{c}\!\sqrt{\frac{180c{+}872}{5c{+}22}}$ \\[4pt]
$\cW_3$
  & $5c/6$ & $\neq 0$ & $\neq 0$ & $\neq 0$
  & $\mathbf{M}$ & $\infty$ & --- \\[2pt]
$\cW_N$ ($N{\geq}3$)
  & $(H_N{-}1)\,c$ & $\neq 0$ & $\neq 0$ & $\neq 0$
  & $\mathbf{M}$ & $\infty$ & --- \\[2pt]
\bottomrule
\end{tabular}
\end{center}

\noindent The ${}^*$ entries for classes C indicate values
on the weight-changing primary line; the quartic contact
invariant lives on a higher-dimensional stratum.
For class~$\mathbf{M}$, the dashes (---) indicate that the
shadow radius depends on the choice of primary line and has
not been computed in closed form for the full multi-channel
problem.


\subsection{Complementarity of discriminants}
\label{subsec:complementarity}

Koszul duality $A \mapsto A^{!}$ transforms the shadow
invariants in a controlled way.  For Virasoro, the chiral
Koszul dual is $\mathrm{Vir}_c^{!} = \mathrm{Vir}_{26-c}$
(the duality $c \mapsto 26 - c$).

\begin{proposition}[Discriminant complementarity]
\label{prop:complementarity}
The critical discriminants satisfy
\begin{equation}\label{eq:complementarity}
  \Delta(\mathrm{Vir}_c) + \Delta(\mathrm{Vir}_{26-c})
  \;=\;
  \frac{6960}{(5c+22)(152-5c)}.
\end{equation}
The numerator $6960 = 40 \times 174$ is independent of~$c$.
At the self-dual point $c = 13$:
$\Delta(13) = 40/87$, and the sum is $80/87$.
The product $(5c+22)(152-5c)$ has maximum $7569$ at $c = 13$.
\end{proposition}

\begin{proof}
Direct computation:
\begin{align*}
  \Delta(c) + \Delta(26-c)
  &\;=\;
  \frac{40}{5c+22} + \frac{40}{5(26-c)+22} \\
  &\;=\;
  \frac{40}{5c+22} + \frac{40}{152-5c} \\
  &\;=\;
  40\cdot\frac{(152-5c)+(5c+22)}{(5c+22)(152-5c)}
  \;=\;
  \frac{40 \times 174}{(5c+22)(152-5c)}.
  \qedhere
\end{align*}
\end{proof}

\begin{proposition}[Shadow radius at special points]
\label{prop:special-radii}
The Virasoro shadow growth rate at selected central
charges:
\begin{center}
\renewcommand{\arraystretch}{1.2}
\begin{tabular}{ccccc}
\toprule
$c$ & $\rho(c)$ & $\rho(26{-}c)$ &
  convergent? & interpretation \\
\midrule
$1/2$ & $12.53$ & $0.224$ & no/yes &
  Ising/$\mathrm{Vir}_{51/2}$ \\
$1$ & $6.242$ & $0.242$ & no/yes &
  free boson/$\mathrm{Vir}_{25}$ \\
$c^*$ & $1$ & $0.430$ & marginal/yes &
  critical \\
$13$ & $0.467$ & $0.467$ & yes/yes &
  self-dual \\
$25$ & $0.242$ & $6.242$ & yes/no &
  dual of $c{=}1$ \\
$26$ & $0.232$ & $\infty$ & yes/--- &
  string/$\kappa^{!}{=}0$ \\
\bottomrule
\end{tabular}
\end{center}
At $c = 13$, the shadow growth rates are equal
($\rho = \rho^{!}$), reflecting the self-duality of
$\mathrm{Vir}_{13}$.  At $c = 26$, the Koszul dual
$\mathrm{Vir}_0$ has $\kappa^{!} = 0$ (the bar complex
is uncurved), and $\rho^{!}$ is undefined.
\end{proposition}

\begin{remark}[Duality and convergence]
\label{rem:duality-convergence}
The Koszul duality $c \leftrightarrow 26-c$ exchanges
convergent and divergent towers: if $\rho(c) < 1$ then
$\rho(26-c) > 1$ generically, and vice versa.  The
self-dual point $c = 13$ is the unique central charge where
$\rho(A) = \rho(A^{!})$.  This reflects the complementarity
principle: the sum
$\Delta(c) + \Delta(26-c) = 6960/[(5c+22)(152-5c)]$
is symmetric in $c \leftrightarrow 26-c$, with maximum at
$c = 13$.
\end{remark}

\section{The \texorpdfstring{$\cW_3$}{W3} cross-channel computation}
\label{sec:genus2}

The genus-$2$ free energy of a single-generator chirally Koszul
algebra is determined by the modular characteristic alone:
$F_2(A) = \kappa(A) \cdot \lambda_2^{\mathrm{FP}}$, where
$\lambda_2^{\mathrm{FP}} = 7/5760$.  For a multi-generator
algebra, the graph sum acquires \emph{cross-channel} contributions
from edges carrying different generators.  The $\cW_3$ algebra
is the simplest instance.

The cross-channel correction, computed by stable graph summation, is
\begin{equation}\label{eq:delta-F2-w3}
  \delta F_2(\cW_3)
  \;=\;
  \frac{c + 204}{16c},
\end{equation}
the first explicit multi-channel genus-$2$ amplitude for any
vertex algebra with generators of distinct conformal weight.

\begin{remark}[Status]
\label{rem:cross-channel-status}
The cross-channel correction~\eqref{eq:delta-F2-w3} uses
per-channel genus-$1$ vertex factors $\kappa_i/24$ without
$R$-matrix corrections.  The open problem
\textup{op:multi-generator-universality} is resolved negatively:
the $R$-matrix dressing does \emph{not} cancel the cross-channel
correction at $g \geq 2$.  The per-channel scalar part
$F_2^{\mathrm{scal}} = (5c/6) \cdot 7/5760 = 7c/6912$ is
unconditional; the cross-channel correction
$\delta F_2 = (c{+}204)/(16c)$ is a genuine additional term.
\end{remark}


\subsection{The $\cW_3$ Frobenius algebra}

The $\cW_3$ algebra has two strong generators:
$T$ (conformal weight~$2$, the Virasoro stress tensor) and
$W$ (conformal weight~$3$, the spin-$3$ current).
The bar-complex CohFT has a two-dimensional Frobenius algebra
with the following data.

\begin{lemma}[$\cW_3$ Frobenius data]
\label{lem:w3-frobenius}
\leavevmode
\begin{enumerate}[(i)]
\item \emph{Curvature metric.}\enspace
  $\eta_{TT} = \kappa_T = c/2$,
  $\eta_{WW} = \kappa_W = c/3$,
  $\eta_{TW} = 0$ \textup{(}diagonal: $T$ and $W$ are
  orthogonal channels\textup{)}.
  The total modular characteristic is
  $\kappa(\cW_3) = \kappa_T + \kappa_W = 5c/6$.

\item \emph{Propagators.}\enspace
  $\eta^{TT} = 2/c$,\enspace $\eta^{WW} = 3/c$.
  Both channels use the weight-$1$ bar propagator
  $d\!\log E(z,w)$, regardless of the conformal weight
  of the generator.

\item \emph{Three-point structure constants.}\enspace
  The $\mathbb{Z}_2$ symmetry $W \mapsto -W$ forces all
  structure constants with an odd number of $W$~indices
  to vanish.  The nonzero constants are
  \[
    C_{TTT} = C_{TWW} = C_{WWT} = c.
  \]
  In particular, $C_{TTW} = C_{WWW} = 0$.

\item \emph{Four-point vertex.}\enspace
  For the genus-$0$ quartic vertex in the $s$-channel
  factorization:
  \[
    V_{0,4}(i,i\,|\,j,j)
    \;=\;
    \sum_{m \in \{T,W\}} \eta^{mm}\,C_{iim}\,C_{jjm}
    \;=\; 2c
  \]
  for all channel pairs $(i,j)$.  The universality follows
  because only the $T$-channel contributes:
  $\eta^{TT} C_{iiT} C_{jjT} = (2/c) \cdot c \cdot c = 2c$.
\end{enumerate}
\end{lemma}

\begin{proof}
Part~(i): the curvature $\kappa_i$ is the leading OPE pole
coefficient.  From $T_{(3)}T = c/2$ and $W_{(5)}W = c/3$.
Part~(iii): from the bar-complex $r$-matrix (OPE poles
shifted by~$-1$ via $d\!\log$):
$T_{(1)}T = 2T$ gives $C^T_{TT} = 2$, so
$C_{TTT} = \eta_{TT} \cdot 2 = c$;
$T_{(1)}W = 3W$ gives $C^W_{TW} = 3$, so
$C_{TWW} = \eta_{WW} \cdot 3 = c$;
$W_{(3)}W = 2T$ gives $C^T_{WW} = 2$, so
$C_{WWT} = \eta_{TT} \cdot 2 = c$.
Part~(iv): direct computation.
\end{proof}


\subsection{Stable graphs at genus $2$}

There are seven stable graphs of genus~$2$ with no marked
points:

\begin{center}
\renewcommand{\arraystretch}{1.2}
\small
\begin{tabular}{lllcc}
\toprule
Graph & Vertices & Edges & $|\Aut|$ & $h^1$ \\
\midrule
$\Gamma_0$: smooth &
  $1 \times (g{=}2, \text{val}=0)$ &
  $0$ & $1$ & $0$ \\
$\Gamma_1$: figure-eight &
  $1 \times (g{=}1, \text{val}=2)$ &
  $1$ self-loop & $2$ & $1$ \\
$\Gamma_2$: banana &
  $1 \times (g{=}0, \text{val}=4)$ &
  $2$ self-loops & $8$ & $2$ \\
$\Gamma_3$: dumbbell &
  $(g{=}1,1) + (g{=}1,1)$ &
  $1$ bridge & $2$ & $0$ \\
$\Gamma_4$: theta &
  $(g{=}0,3) + (g{=}0,3)$ &
  $3$ bridges & $12$ & $2$ \\
$\Gamma_5$: lollipop &
  $(g{=}0,3) + (g{=}1,1)$ &
  $1$ self-loop $+$ $1$ bridge & $2$ & $1$ \\
$\Gamma_6$: barbell &
  $(g{=}0,3) + (g{=}0,3)$ &
  $2$ self-loops $+$ $1$ bridge & $8$ & $2$ \\
\bottomrule
\end{tabular}
\end{center}

\noindent Each graph satisfies the genus formula
$g = h^1 + \sum_v g_v = 2$ and the stability condition
$2g_v + \text{val}_v \geq 3$ at every vertex.

Each edge carries a channel label $\sigma(e) \in \{T, W\}$.
The amplitude of a graph $\Gamma$ with channel assignment
$\sigma\colon E(\Gamma) \to \{T, W\}$ is
\begin{equation}\label{eq:feynman-rule}
  A(\Gamma, \sigma)
  \;=\;
  \frac{1}{|\Aut(\Gamma)|}
  \prod_{e \in E(\Gamma)} \eta^{\sigma(e),\sigma(e)}
  \prod_{v \in V(\Gamma)} V_v(g_v,\, \text{channels}_v),
\end{equation}
where $V_v$ is the vertex factor: $C_{ijk}$ for genus-$0$
trivalent vertices, $V_{0,4}$ for genus-$0$ four-valent
vertices, and $\kappa_i/24$ for genus-$1$ vertices with
a single half-edge of channel~$i$.

An assignment $\sigma$ is \emph{mixed} if not all edges
carry the same channel.  The cross-channel correction is
the sum of mixed-assignment amplitudes over all graphs.


\subsection{Graph-by-graph computation}

We compute the mixed-channel contribution from each
of the six boundary graphs ($\Gamma_1$--$\Gamma_6$;
the smooth graph $\Gamma_0$ carries no edges).

\subsubsection*{$\Gamma_1$: figure-eight}

One edge, one self-loop at the genus-$1$ vertex.
The single edge carries either $T$ or~$W$; there is no
mixed assignment.  Cross-channel contribution: $0$.

\subsubsection*{$\Gamma_2$: banana}

Two self-loops at a genus-$0$ four-valent vertex, with
$|\Aut| = 8$.
The four channel assignments $(\sigma_1, \sigma_2) \in
\{T,W\}^2$ and the quartic vertex universality
$V_{0,4}(i,i\,|\,j,j) = 2c$ give:
\begin{align*}
  (T,T):&\quad \eta^{TT}\,\eta^{TT}\cdot 2c
    = (2/c)^2 \cdot 2c = 8/c, \\
  (W,W):&\quad \eta^{WW}\,\eta^{WW}\cdot 2c
    = (3/c)^2 \cdot 2c = 18/c, \\
  (T,W)\text{ or }(W,T):&\quad
    \eta^{TT}\,\eta^{WW}\cdot 2c
    = (2/c)(3/c) \cdot 2c = 12/c,
    \quad\text{two assignments.}
\end{align*}
The mixed raw amplitude is $2 \cdot 12/c = 24/c$.
Dividing by $|\Aut| = 8$:
\begin{equation}\label{eq:banana-cross}
  \delta_{\mathrm{banana}}
  \;=\; \frac{24}{8c}
  \;=\; \frac{3}{c}.
\end{equation}

\subsubsection*{$\Gamma_3$: dumbbell}

One bridge connecting two genus-$1$ vertices.  The single
edge carries either $T$ or $W$; the off-diagonal metric
$\eta^{TW} = 0$ kills any notion of mixing.
Cross-channel contribution: $0$.

\subsubsection*{$\Gamma_4$: theta}

Three bridges connecting two genus-$0$ trivalent vertices,
with $|\Aut| = 12$.
There are $2^3 = 8$ channel assignments.

The $\mathbb{Z}_2$ parity $C_{TTW} = C_{WWW} = 0$ kills
every assignment in which either vertex receives an odd number
of $W$ half-edges.  Since each vertex sees exactly one
half-edge from each bridge, the constraint is that the number
of $W$-labeled edges must be even.

\begin{itemize}
\item $(T,T,T)$: all-$T$, diagonal.
  Amplitude: $(2/c)^3 \cdot c^2 = 8/c$.  With $1/12$:
  $2/(3c)$.

\item $(T,T,W)$ and permutations (3~assignments):
  each vertex gets $(T,T,W)$.
  $C_{TTW} = 0$, so all three vanish.

\item $(T,W,W)$ and permutations (3~assignments):
  each vertex gets $(T,W,W)$.
  $C_{TWW} = c$ at both vertices.
  Propagators: $\eta^{TT}\,(\eta^{WW})^2
  = (2/c)(3/c)^2 = 18/c^3$.
  Raw: $18 c^2/c^3 = 18/c$.
  Three such assignments; total mixed raw: $54/c$.
  With $1/12$: $54/(12c) = 9/(2c)$.

\item $(W,W,W)$: $C_{WWW} = 0$, vanishes.
\end{itemize}

\noindent The mixed contribution from the theta graph is
\begin{equation}\label{eq:theta-cross}
  \delta_{\mathrm{theta}}
  \;=\; \frac{54}{12c}
  \;=\; \frac{9}{2c}.
\end{equation}

\subsubsection*{$\Gamma_5$: lollipop}

One self-loop at the genus-$0$ trivalent vertex~$v_0$,
one bridge from~$v_0$ to the genus-$1$ vertex~$v_1$.
$|\Aut| = 2$.
Write $(\sigma_1, \sigma_2)$ for the channels of the
self-loop and the bridge.

\begin{itemize}
\item $(T,W)$: self-loop $T$, bridge~$W$.
  $v_0$ gets half-edges $(T,T,W)$:
  $C_{TTW} = 0$.  Vanishes.

\item $(W,T)$: self-loop $W$, bridge~$T$.
  $v_0$ gets half-edges $(W,W,T)$:
  $C_{WWT} = c$.
  $v_1$ gets half-edge $T$:
  $V_{1,1}(T) = \kappa_T/24 = c/48$.
  Propagators:
  $\eta^{WW}\,\eta^{TT} = (3/c)(2/c) = 6/c^2$.
  Raw: $(6/c^2) \cdot c \cdot (c/48) = 1/8$.
  With $1/|\Aut| = 1/2$:
  $1/16$.
\end{itemize}

\noindent The lollipop mixed amplitude is
\begin{equation}\label{eq:lollipop-cross}
  \delta_{\mathrm{lollipop}}
  \;=\; \frac{1}{16}.
\end{equation}
This is \emph{independent of the central charge~$c$}:
$(3/c)(2/c) \cdot c \cdot (c/48) = 1/8$, and $1/8$
divided by $|\Aut| = 2$ is~$1/16$.

\subsubsection*{$\Gamma_6$: barbell}

Two genus-$0$ trivalent vertices~$v_1, v_2$, each carrying a
self-loop, connected by a single bridge.
$|\Aut| = 8$ (swapping $v_1 \leftrightarrow v_2$ and
independently swapping the two half-edges of each self-loop).

Write $(\sigma_1, \sigma_2, \sigma_3)$ for the channels of
the self-loop at~$v_1$, the self-loop at~$v_2$, and the bridge.
Vertex~$v_i$ receives half-edges $(\sigma_i, \sigma_i, \sigma_3)$
and contributes the three-point function~$C_{\sigma_i \sigma_i \sigma_3}$.

\begin{itemize}
\item $(T,W,T)$: $v_1$ gets $(T,T,T)$,
  $C_{TTT} = c$;\enspace $v_2$ gets $(W,W,T)$,
  $C_{WWT} = c$.
  Propagators: $\eta^{TT}\,\eta^{WW}\,\eta^{TT}
  = (2/c)(3/c)(2/c) = 12/c^3$.
  Raw: $12 c^2/c^3 = 12/c$.

\item $(W,T,T)$: by the $v_1 \leftrightarrow v_2$ symmetry,
  same raw amplitude $12/c$.

\item $(W,W,T)$: $v_1$ gets $(W,W,T)$,
  $C_{WWT} = c$;\enspace $v_2$ gets $(W,W,T)$,
  $C_{WWT} = c$.
  Propagators: $(\eta^{WW})^2\,\eta^{TT}
  = (3/c)^2(2/c) = 18/c^3$.
  Raw: $18c^2/c^3 = 18/c$.

\item All other mixed assignments vanish:
  any assignment with bridge channel~$W$ requires
  $C_{\sigma_i \sigma_i W}$;\enspace
  $C_{TTW} = 0$ and $C_{WWW} = 0$.
\end{itemize}
Total mixed raw amplitude: $12/c + 12/c + 18/c = 42/c$.
Dividing by $|\Aut| = 8$:
\begin{equation}\label{eq:barbell-cross}
  \delta_{\mathrm{barbell}}
  \;=\; \frac{42}{8c}
  \;=\; \frac{21}{4c}.
\end{equation}


\subsection{The cross-channel correction}

\begin{theorem}[$\cW_3$ cross-channel correction]
\label{thm:genus2}
The cross-channel correction to the genus-$2$ free energy
of~$\cW_3$ is
\begin{equation}\label{eq:cross-channel-total}
  \delta F_2(\cW_3)
  \;=\;
  \underbrace{\frac{3}{c}}_{\text{banana}}
  \;+\;
  \underbrace{\frac{9}{2c}}_{\text{theta}}
  \;+\;
  \underbrace{\frac{1}{16}}_{\text{lollipop}}
  \;+\;
  \underbrace{\frac{21}{4c}}_{\text{barbell}}
  \;=\;
  \frac{c + 204}{16c}.
\end{equation}
The open problem \textup{op:multi-generator-universality}
is resolved negatively.  The full genus-$2$
free energy decomposes as
\[
  F_2(\cW_3)
  \;=\;
  \underbrace{\frac{5c}{6}\cdot\frac{7}{5760}}_{\displaystyle
    F_2^{\mathrm{scal}} = 7c/6912}
  \;\;+\;\;
  \underbrace{\frac{c + 204}{16c}}_{\displaystyle
    \delta F_2^{\mathrm{cross}}}.
\]
\end{theorem}

\begin{proof}
The cross-channel contributions from all seven graphs
are computed in
\eqref{eq:banana-cross}--\eqref{eq:barbell-cross}.
Four graphs contribute: banana ($3/c$), theta ($9/(2c)$),
lollipop ($1/16$), barbell ($21/(4c)$).  The remaining three (smooth, figure-eight,
dumbbell) have zero mixed amplitude: the smooth graph has no edges,
while the figure-eight and dumbbell each have a single edge, so
every assignment is diagonal.

For the sum:
\[
  \frac{3}{c} + \frac{9}{2c} + \frac{1}{16} + \frac{21}{4c}
  \;=\;
  \frac{48}{16c} + \frac{72}{16c} + \frac{c}{16c} + \frac{84}{16c}
  \;=\;
  \frac{c + 204}{16c}.
  \qedhere
\]
\end{proof}

\begin{remark}[Properties of the cross-channel correction]
\label{rem:cross-channel-properties}
\leavevmode
\begin{enumerate}[(i)]
\item \emph{The $c$-independent piece.}\enspace
  The lollipop contribution $1/16$ is the only
  $c$-independent term:
  $\eta^{WW}\,\eta^{TT}\,C_{WWT}\,({\kappa_T}/{24})
  = (3/c)(2/c) \cdot c \cdot (c/48) = 1/8$,
  and every factor of~$c$ cancels.  As $c \to \infty$,
  $\delta F_2 \to 1/16$.

\item \emph{The rational piece.}\enspace
  The $c$-dependent terms total $51/(4c)$: banana $3/c$,
  theta $9/(2c)$, and barbell $21/(4c)$.  All three arise
  from genus-$0$ vertices, mediated by the $\mathbb{Z}_2$-allowed
  mixed assignments.  They dominate at small~$c$.

\item \emph{Vanishing at $c = -204$.}\enspace
  The cross-channel correction vanishes at $c = -204$.

\item \emph{Relative magnitude.}\enspace
  The ratio
  $\delta F_2 / F_2^{\mathrm{scal}}
  = 432(c+204)/(7c^2)$
  tends to zero as $c \to \infty$: the cross-channel
  correction is a subleading perturbation of the scalar
  part for large central charge.
  For example, at $c = 50$:
  \[
    F_2^{\mathrm{scal}}(50) = \frac{7 \cdot 50}{6912}
    = \frac{175}{3456},
    \qquad
    \delta F_2(50) = \frac{254}{800}
    = \frac{127}{400}.
  \]
\end{enumerate}
\end{remark}

\begin{remark}[The $\mathbb{Z}_2$ selection rule]
\label{rem:z2-selection}
The $\mathbb{Z}_2$ symmetry $W \mapsto -W$ of the $\cW_3$
algebra forces $C_{TTW} = C_{WWW} = 0$, eliminating every
mixed assignment in which a vertex receives an odd number of
$W$-legs.  For the theta graph, this kills the three
$(T,T,W)$-type assignments and preserves the three
$(T,W,W)$-type assignments.  For $\cW_N$ with $N \geq 4$,
the additional generators (spins $4, \ldots, N$) break this
symmetry.
\end{remark}

\section{The shadow connection and Koszul duality}
\label{sec:connection}

\begin{theorem}[Shadow connection]
\label{thm:connection}
The shadow metric $Q(t)$ defines a logarithmic connection
\[
  \nabla^{\mathrm{sh}}
  \;=\; d \;-\; \frac{Q'(t)}{2\,Q(t)}\,dt
\]
on $L \setminus \{Q = 0\}$.
\begin{enumerate}[(i)]
\item The connection has regular singularities of
  residue~$\tfrac{1}{2}$ at each zero of~$Q$.
\item The monodromy around each branch point is~$-1$.
\item The flat section with $\Phi(0) = 1$ is
  $\Phi(t) = \sqrt{Q(t)/Q(0)}$, and
  $H(t) = 2\kappa\,t^2\,\Phi(t)$.
\item Under $A \mapsto A^{!}$ (Koszul duality), for
  Virasoro:
  \[
    \Delta(\mathrm{Vir}_c) + \Delta(\mathrm{Vir}_{26-c})
    \;=\; \frac{6960}{(5c+22)(152-5c)}.
  \]
\end{enumerate}
\end{theorem}

\begin{proof}
Parts (i)--(iii) are proved in
Proposition~\ref{prop:conn-properties}.
Part~(iv) is Proposition~\ref{prop:complementarity}.
\end{proof}

\section{Outlook}\label{sec:outlook}

Three directions extend the algebraicity theorem.


\subsection*{Shadow obstruction towers for non-standard families}

The standard landscape (Heisenberg, affine Kac--Moody,
$\beta\gamma$, Virasoro, $\cW_N$) exhausts the four shadow
depth classes G, L, C, M.  Non-standard families
(admissible-level quotients, $N=2$ superconformal algebras,
non-simply-laced affine vertex algebras at fractional level)
pose new questions.  For admissible-level quotients
$L_k(\fg)$, null vectors modify the bar complex and may
shift the shadow depth class.  The Bershadsky--Polyakov
algebra $\cW_3^{(2)}$ at admissible levels and the triplet
$\cW(p)$ algebras are concrete test cases.  Whether the
four-class partition exhausts all chirally Koszul algebras,
or whether exotic depth classes (e.g., $r_{\max} = 5$) arise
outside the standard landscape, is open.


\subsection*{The Pixton connection}

For class~$\mathbf{M}$ algebras (Virasoro, $\cW_N$), the
infinite shadow obstruction tower $\{S_r\}_{r \geq 2}$ produces an infinite
sequence of tautological classes on $\Mbar_{g,n}$ via the
Maurer--Cartan equation.  At genus~$2$--$3$, computations in
the modular convolution algebra yield tautological relations
that belong to the Pixton ideal $\mathcal{I}_{\mathrm{Pix}}
\subset R^*(\Mbar_{g,n})$, the ideal of tautological
relations constructed by Pixton and proved to hold
by Janda--Pandharipande--Pixton--Zvonkine \cite{JPPZ}
(completeness is conjectural in general).
The MC equation $D^2 = 0$ forces algebraic relations
among stable-graph amplitudes at each genus, and these
project to tautological identities on $\Mbar_g$.

The conjecture is that the Virasoro MC tower generates the
entire Pixton ideal, via the shadow-to-strata map sending
each MC obstruction $o_r$ to a codimension-$(r-2)$
tautological class.  This is verified at genus~$2$ (the
planted-forest formula
$\delta_{\mathrm{pf}}^{(2,0)} = S_3(10S_3 - \kappa)/48$
lies in the strata algebra) and at genus~$3$ (the
$11$-term polynomial in $\kappa, S_3, S_4, S_5$ matches
known relations).  The main obstruction is the flat unit
question for the shadow CohFT, needed to apply
Givental--Teleman reconstruction.


\subsection*{Analytic completion and shadow convergence}

The shadow growth rate $\rho(A)$
(Definition~\ref{def:shadow-radius}) divides the Virasoro
family into a convergent regime ($c > c^*$, $\rho < 1$) and
a divergent regime ($c < c^*$, $\rho > 1$).  The convergent
shadow obstruction tower defines a shadow partition function
$Z^{\mathrm{sh}}(q, t)$ with absolute convergence in both the
genus variable~$q$ and the arity variable~$t$.  The
string partition function
$Z^{\mathrm{str}}(g_s) = \sum_g F_g\,g_s^{2g-2}$
diverges factorially as $(2g)!$ by Shenker's argument.

The shadow obstruction tower converges because its coefficients decay
as $S_r \sim \rho^r\,r^{-5/2}$ (algebraic, from a
square-root singularity), while string amplitudes grow as
$(2g)!$ (factorial, from the volume of $\Mbar_g$).  The
shadow partition function is therefore a resummation of the
factorially divergent string expansion, converging on a
domain determined by the shadow radius.

The passage from algebraic shadow obstruction tower to analytic
partition function requires the sewing envelope
construction: a Hausdorff completion of the vertex algebra
in the all-amplitude seminorm topology.  The
Hilbert--Schmidt sewing condition (that the sewing
operator has finite Hilbert--Schmidt norm) is satisfied for
the entire standard landscape and implies trace-class
closed amplitudes at all genera.  The resulting analytic bar
coalgebra lives in the coderived category at genus $g \geq 1$,
where curvature obstructs the ordinary derived category.

%%% ================================================================
%%%  BIBLIOGRAPHY
%%% ================================================================

\begin{thebibliography}{99}

\bibitem{ADKMV}
M.~Aganagic, R.~Dijkgraaf, A.~Klemm, M.~Mari\~no, and C.~Vafa,
\emph{Topological strings and integrable hierarchies},
Comm.\ Math.\ Phys.\ \textbf{261} (2006), 451--516.

\bibitem{BD}
A.~Beilinson and V.~Drinfeld,
\emph{Chiral Algebras},
AMS Colloquium Publications, vol.~51, 2004.

\bibitem{CG}
K.~Costello and O.~Gwilliam,
\emph{Factorization Algebras in Quantum Field Theory},
vols.~1--2, Cambridge University Press, 2017/2021.

\bibitem{DijkgraafVafa}
R.~Dijkgraaf and C.~Vafa,
\emph{Matrix models, topological strings, and supersymmetric
gauge theories},
Nuclear Phys.\ B \textbf{644} (2002), 3--20.

\bibitem{EO}
B.~Eynard and N.~Orantin,
\emph{Invariants of algebraic curves and topological expansion},
Comm.\ Number Theory Phys.\ \textbf{1} (2007), 347--452.

\bibitem{FP}
C.~Faber and R.~Pandharipande,
\emph{Hodge integrals and moduli spaces of curves},
in:\ Handbook of Moduli, vol.~I, International Press, 2013.

\bibitem{FrenkelBenZvi}
E.~Frenkel and D.~Ben-Zvi,
\emph{Vertex Algebras and Algebraic Curves},
2nd ed., Mathematical Surveys and Monographs, vol.~88,
AMS, 2004.

\bibitem{GK}
E.~Getzler and M.~M.~Kapranov,
\emph{Modular operads},
Compositio Math.\ \textbf{110} (1998), 65--126.

\bibitem{JPPZ}
F.~Janda, R.~Pandharipande, A.~Pixton, and D.~Zvonkine,
\emph{Double ramification cycles on the moduli spaces of curves},
Publ.\ Math.\ Inst.\ Hautes \'Etudes Sci.\ \textbf{125} (2017), 221--266.

\bibitem{Kac-vertex}
V.~G.~Kac,
\emph{Vertex Algebras for Beginners},
2nd ed., University Lecture Series, vol.~10, AMS, 1998.

\bibitem{KontsevichSoibelman}
M.~Kontsevich and Y.~Soibelman,
\emph{Stability structures, motivic Donaldson--Thomas
invariants and cluster transformations},
arXiv:0811.2435.

\bibitem{LV}
J.-L.~Loday and B.~Vallette,
\emph{Algebraic Operads},
Grundlehren der Mathematischen Wissenschaften, vol.~346,
Springer, 2012.

\bibitem{Zhu}
Y.~Zhu,
\emph{Modular invariance of characters of vertex operator
algebras},
J.~Amer.\ Math.\ Soc.\ \textbf{9} (1996), 237--302.

\end{thebibliography}

\end{document}
